\documentclass[11pt]{article}

\usepackage[margin=1in]{geometry}
\usepackage[T1]{fontenc}
\usepackage[utf8]{inputenc}
\usepackage{lmodern}
\usepackage{microtype}
\usepackage{amsmath,amssymb,amsthm,mathtools,bm}
\usepackage{booktabs,tabularx,array,multirow}
\usepackage{enumitem}
\usepackage{algorithm}
\usepackage[noend]{algpseudocode}
\makeatletter
\renewcommand{\theHALG@line}{\thealgorithm.\arabic{ALG@line}}
\makeatother
\usepackage{xcolor}
\usepackage[hidelinks]{hyperref}
\usepackage[nameinlink,capitalize,noabbrev]{cleveref}
\usepackage{natbib}
\hypersetup{
  pdftitle={Primary Ties and the Price of Unknown Trade-offs in Lexicographic Bandits},
  pdfauthor={Colombo Team},
  pdfsubject={Simultaneous regret, trade-off envelopes, and metric complexity},
  pdfkeywords={lexicographic bandits, minimax regret, adaptation, Lipschitz bandits}
}
\emergencystretch=1.5em

\allowdisplaybreaks
\setlength{\parindent}{1.2em}
\setlength{\parskip}{0.15em}
\setlist{nosep}

\newtheorem{theorem}{Theorem}[section]
\newtheorem{lemma}[theorem]{Lemma}
\newtheorem{proposition}[theorem]{Proposition}
\newtheorem{corollary}[theorem]{Corollary}
\newtheorem{definition}[theorem]{Definition}
\newtheorem{assumption}[theorem]{Assumption}
\newtheorem{remark}[theorem]{Remark}
\theoremstyle{definition}
\newtheorem{example}[theorem]{Example}

\newcommand{\X}{\mathcal X}
\newcommand{\A}{\mathcal A}
\newcommand{\C}{\mathcal C}
\newcommand{\F}{\mathcal F}
\newcommand{\Pcal}{\mathcal P}
\newcommand{\Mcal}{\mathcal M}
\newcommand{\Ocal}{\mathcal O}
\newcommand{\E}{\mathbb E}
\newcommand{\Pp}{\mathbb P}
\newcommand{\R}{\mathbb R}
\newcommand{\N}{\mathbb N}
\newcommand{\KL}{\operatorname{KL}}
\newcommand{\kl}{\operatorname{kl}}
\newcommand{\Pack}{\operatorname{Pack}}
\newcommand{\Cov}{\operatorname{Cov}}
\newcommand{\eps}{\varepsilon}
\newcommand{\ind}{\mathbf 1}
\newcommand{\pos}[1]{[#1]_+}
\newcommand{\wtO}{\widetilde O}
\newcommand{\wtOmega}{\widetilde\Omega}
\newcommand{\wtTheta}{\widetilde\Theta}
\newcommand{\LexETC}{{\normalfont\textsc{Lex-ETC}}}
\newcommand{\PFETC}{{\normalfont\textsc{PF-ETC}}}
\newcommand{\RSETC}{{\normalfont\textsc{RS-ETC}}}
\newcommand{\argmax}{\operatorname*{arg\,max}}
\newcommand{\argmin}{\operatorname*{arg\,min}}

\newcolumntype{Y}{>{\raggedright\arraybackslash}X}
\crefname{assumption}{Assumption}{Assumptions}
\Crefname{assumption}{Assumption}{Assumptions}
\crefname{definition}{Definition}{Definitions}
\Crefname{definition}{Definition}{Definitions}
\crefname{remark}{Remark}{Remarks}
\Crefname{remark}{Remark}{Remarks}

\title{Primary Ties and the Price of Unknown Trade-offs\\
       in Lexicographic Bandits}
\author{Colombo Team\\[0.2em]
  \normalsize\href{mailto:colombo.2026@outlook.com}{\texttt{colombo.2026@outlook.com}}}
\date{}

\begin{document}
\maketitle

\begin{abstract}
Lexicographic bandits rank objectives by priority: the secondary objective
is optimized only among primary-optimal arms. Without side information,
an exact primary tie is statistically difficult to distinguish from a small
primary improvement. Building on the classical two-arm obstruction, we
characterize its consequences for arm count, trade-off information, and
metric complexity. For two objectives and $K$ arms, the simultaneous
expected minimax regret is
$\widetilde\Theta(\min\{T,K^{1/3}T^{2/3}\})$.
When the learner is given an upper bound $L$ on the local trade-off,
we prove the phase diagram
\[
  \widetilde\Theta\!\left(
  \min\{T,\sqrt{(1\vee L)KT},K^{1/3}T^{2/3}\}\right).
\]
The upper bounds control the positive part of each coordinate regret,
whereas the lower bounds already hold for general coordinate regret.
In compact metric spaces, a multiscale filtering algorithm adapts to the
packing profile of the primary near-optimal set without knowing that
profile. On regular anchor-plus-tie spaces with packing growth $Zr^{-d}$,
the resulting rate is
$\widetilde\Theta_d(\min\{T,Z^{1/(d+3)}T^{(d+2)/(d+3)}\})$.
We also give a critical-entropy characterization for arbitrary compact tie
spaces, an explicit expected-regret Pareto lower constraint, and the sharp
two-arm confidence dependence
$\Theta(\min\{T,T^{2/3}\log^{1/3}(1/\delta)\})$.
The metric guarantees are information-theoretic and assume access to
maximal nets. The analysis separates primary-tie complexity from the
positive-gap annular geometry used in supplied-trade-off guarantees.
\end{abstract}

\medskip
\noindent\textbf{Keywords:} lexicographic bandits; simultaneous regret;
minimax bounds; unknown trade-offs; Lipschitz bandits; near-optimality packing.
\par\smallskip
\begin{samepage}
\noindent\textbf{AI Usage.}
Nearly the entire research pipeline for this paper was carried out by
Colombo, powered by GPT-5.6 Sol.
\par
\end{samepage}

\section{Introduction}
\label{sec:introduction}

In a lexicographic decision problem, objectives have a strict priority
order. A learner first seeks the best primary reward and then seeks the
best secondary reward among primary-optimal actions. This ordering differs
from optimizing a fixed weighted sum: even a small primary improvement
takes precedence over a large secondary loss. With noisy observations,
the central difficulty is deciding whether two primary means are exactly
equal or merely close.

Consider two arms whose primary means coincide. One arm has a fixed
secondary advantage, so sampling the other arm incurs secondary regret.
If the latter arm instead has a small, unknown primary advantage, repeatedly
choosing the former incurs primary regret. \citet{huyuk2021lexicographic} used this obstruction to prove a
$T^{2/3}$ lower bound for prior-free lexicographic bandits. Their example
also shows why identifying the primary-optimal set is more demanding than
approximately maximizing a single reward function.

Subsequent work introduced a local trade-off parameter $\lambda_\star$
that bounds the secondary improvement available per unit of primary loss
\citep{xue2024lipschitz,xue2025jmlr}. Supplying a valid upper bound on
$\lambda_\star$ permits stronger coordinate-wise regret guarantees.
This raises a more quantitative question: how does the statistical cost
depend on the amount of trade-off information and on the size and geometry
of the region that may contain a primary optimum?

This paper answers that question for two objectives. We require one policy
to control both coordinate regrets on every admissible instance. For a
finite set of $K$ arms, the minimax rate without an envelope is
$\widetilde\Theta(K^{1/3}T^{2/3})$, capped by $T$. If an envelope $L$ is
supplied, the rate improves to $\widetilde\Theta(\sqrt{(1\vee L)KT})$
until $L$ reaches order $(T/K)^{1/3}$, at which point the prior-free rate
is already optimal. Thus knowing a very loose envelope has no minimax
benefit at the given horizon.

The same calculation identifies the relevant geometry in a metric arm
space. Suppose there are $N(\eps)$ separated locations where a primary
improvement of height $\eps$ could be hidden. Distinguishing one such
location requires order $\eps^{-2}$ observations. On the null instance,
searching these locations incurs a fixed secondary loss per observation;
on an alternative, missing the improvement incurs order $T\eps$ primary
regret. The balance is therefore
\begin{equation}
\label{eq:intro-balance}
  N(\eps)\eps^{-2}\asymp T\eps,
  \qquad\text{or equivalently}\qquad
  T\eps^3\asymp N(\eps).
\end{equation}
When $N(\eps)\asymp Z\eps^{-d}$, this balance yields the denominator
$d+3$. We turn it into matching minimax bounds on regular metric spaces
and a one-scale entropy formula without a power-law assumption.

\paragraph{Contributions and relation to the classical obstruction.}
The two-arm construction is included as a self-contained starting point
and is explicitly attributed to \citet{huyuk2021lexicographic}. Our
refinements concern the following questions.
\begin{enumerate}[leftmargin=2.1em]
  \item \textbf{Dependence on the number of arms.}
  We match $\Omega(\min\{T,K^{1/3}T^{2/3}\})$ with a simple
  explore-filter-commit upper bound, up to a logarithmic factor.
  The upper bound controls positive coordinate losses even when general
  secondary regret can be negative.
  \item \textbf{Value of a supplied trade-off envelope.}
  We characterize the complete finite-arm transition between square-root
  regret and the cubic-search rate. Robust-score elimination reduces
  exploration losses, while secondary maximization at commitment avoids
  multiplying the final accuracy by the envelope.
  \item \textbf{Adaptation to primary near-optimality geometry.}
  A multiscale algorithm uses neither a trade-off envelope nor a packing
  dimension or coefficient. We prove an instance-wise oracle upper bound,
  matching class-wise rates on regular anchor-plus-tie spaces, and a
  metric-wise critical-entropy characterization.
  \item \textbf{Expected-regret and confidence refinements.}
  We derive an explicit constraint between null-secondary and
  alternative-primary expected regret. For high-probability regret,
  the dependence on $\log(1/\delta)$ is sharp on two arms; for general
  $K$, the upper and lower bounds differ only in an additional
  $\log^{1/3}K$ factor.
\end{enumerate}

The packing profiles, maximal nets, and change-of-measure inequalities
used below are standard. The analysis combines them with the
lexicographic testing obstruction to quantify how side information and
primary near-optimality geometry determine simultaneous regret.
Throughout, a matching metric result refers to the stated primary-tie
class; it does not identify the primary-tie dimension with an
objective-wise annular zooming dimension.

\paragraph{Organization.}
We first place the results relative to the closest bandit guarantees and
specify the model. Sections~\ref{sec:two-arm}--\ref{sec:envelope} develop
the two-arm obstruction and the finite-arm phase diagram.
Section~\ref{sec:profile} treats metric spaces, and
Section~\ref{sec:frontiers} gives the Pareto and confidence refinements.
The appendices contain the remaining proofs.

\section{Related Work}
\label{sec:related}

\paragraph{Finite-arm lexicographic learning.}
\citet{huyuk2021lexicographic} introduced prior-informed and prior-free
lexicographic bandit algorithms. For a fixed number of arms, their
PF-LEX algorithm has a $\widetilde O(T^{2/3})$ guarantee under
priority-based regret, and their Theorem~2 proves a matching
$\Omega(T^{2/3})$ lower bound on two arms. In priority-based regret,
secondary loss is counted only when the chosen arm is primary-optimal.
Our general coordinate regret counts secondary loss on every round;
our upper bounds also control its positive part. On the two-arm lower
family, the relevant coordinate regrets agree under these conventions.
The earlier lower bound therefore applies directly. We extend its
quantitative analysis rather than claim a new $T^{2/3}$ obstruction.

\paragraph{Trade-off information and metric arms.}
\citet{xue2024lipschitz,xue2025jmlr} study general coordinate regret
with a supplied local-trade-off envelope. Their ADLO guarantee uses
objective-wise positive-gap annular dimensions $d_z^i$ and coefficients
$Z_i$. Their two-objective UCB-ADLO method removes the trade-off input
but uses $d_z^1$ and $Z_1$. Table~\ref{tab:related} records the reported
bounds with these inputs separated. Our dimension $d$ instead describes
primary near-optimality or primary-tie packing; the two dimensions are
not interchangeable. More precisely, their annular profile covers
\[
 \X_L^i(r)=\{x:\Lambda_i(L)r/2<\Delta_i(x)\le\Lambda_i(L)r\}
\]
by balls of radius $r/96$, with covering number at most
$Z_i r^{-d_z^i}$. This excludes zero primary gaps, whereas our
primary near-optimality profile includes them.

For finite arms, the supplied-envelope ADLO bound is
\begin{equation}
\label{eq:adlo-finite-comparison}
 (R_1(T),R_2(T))
 \le \widetilde O\!\left(\sqrt{KT},(1+L)\sqrt{KT}\right).
\end{equation}
Our $\sqrt{(1\vee L)KT}$ branch optimizes the maximum of the two
expected regrets. It does not dominate the primary guarantee in
\eqref{eq:adlo-finite-comparison}. The annular-profile lower bound in
\citet[Theorem~3]{xue2025jmlr} has a different quantifier structure:
its benchmark depends on the constructed instance's own annular
profile. It is not a fixed-$(K,L)$ minimax lower bound obtained by
replacing those covering numbers with $K$.

\begin{table}[t]
\centering
\small
\setlength{\tabcolsep}{5pt}
\renewcommand{\arraystretch}{1.16}
\caption{Closest guarantees and the information they use.
All rates suppress logarithms and universal constants.
For two objectives, $\Lambda_1(u)=1$ and $\Lambda_2(u)=1+u$.
The annular parameters $(d_z^i,Z_i)$ and the primary-tie parameters
$(d,Z)$ describe different sets. The notation $\mathfrak M$ denotes
the simultaneous expected minimax criterion defined in
\eqref{eq:joint-risk}; it is not an individual-coordinate optimum.}
\label{tab:related}
\begin{tabularx}{\textwidth}{@{}>{\raggedright\arraybackslash}p{3.0cm}>{\raggedright\arraybackslash}p{3.25cm}Y@{}}
\toprule
Work and arm model & Information / criterion & Guarantee \\
\midrule
H\"uy\"uk--Tekin (2021), finite arms
& No prior reward information; priority-based regret
& $\widetilde O(T^{2/3})$ for fixed $K$;
$\Omega(T^{2/3})$ already for $K=2$. \\
\addlinespace
Xue et al. (2025a), ADLO, metric arms
& Supplied $L\ge\lambda_\star$; general coordinate regret
& $\widetilde O\!\left(\Lambda_i(L)
Z_i^{1/(d_z^i+2)}T^{(d_z^i+1)/(d_z^i+2)}\right)$. \\
\addlinespace
Xue et al. (2025a), UCB-ADLO, metric arms
& Two objectives; supplied $(d_z^1,Z_1)$; no trade-off input
& Reported $\widetilde O\!\left(T^{2/3}+
\Lambda_i(\lambda_\star)Z_i^{1/(d_z^i+2)}
T^{(3d_z^i+4)/(3d_z^i+6)}\right)$, for $\lambda_\star\ge1$. \\
\addlinespace
This paper, finite arms
& No trade-off or optimal-set information; joint criterion
& $\mathfrak M_{T,K}=
\widetilde\Theta(\min\{T,K^{1/3}T^{2/3}\})$. \\
\addlinespace
This paper, finite arms
& Supplied $L\ge\lambda_\star$; joint criterion
& $\mathfrak M_{T,K}(L)=\widetilde\Theta\!\left(
\min\{T,\sqrt{(1\vee L)KT},K^{1/3}T^{2/3}\}\right)$. \\
\addlinespace
This paper, regular tie spaces
& Maximal-net access; no trade-off or profile input
& $\widetilde\Theta_d\!\left(
\min\{T,Z^{1/(d+3)}T^{(d+2)/(d+3)}\}\right)$;
fixed regularity constants. \\
\bottomrule
\end{tabularx}
\end{table}

\paragraph{Other lexicographic models.}
\citet{xue2025linear} study linear bandits with multiple supplied
trade-offs. \citet{xue2025corruptions} give problem-dependent bounds
under adversarial corruptions; both their known- and unknown-corruption
budget algorithms still use a valid trade-off envelope.
\citet{xue2026lexelim} combine regret minimization with best-arm
identification: their outer-layer method uses optimal-set cardinalities,
whereas their inner-layer method uses a trade-off envelope.
Dominant-objective contextual bandits, similarity-based contextual
bandits, and lexicographic reinforcement learning use additional
structure or different performance criteria
\citep{tekin2018dominant,turgay2018similarity,skalse2022lexicographic}.

\paragraph{Metric complexity and adaptation.}
Packing profiles and adaptive discretization are standard in metric
and $X$-armed bandits
\citep{kleinberg2008metric,bubeck2011xarmed,slivkins2011implicit}.
Problem-dependent Lipschitz bounds are developed by
\citet{magureanu2014lipschitz}. We use these geometric tools to measure
the region in which a hidden primary improvement must be ruled out.
Polynomial costs of adaptation also occur when the unknown quantity
is smoothness \citep{locatelli2018adaptivity,hadiji2019polynomial}.
Here smoothness is fixed; the unknown trade-off concerns the relation
between two rewards. Our lower bounds use standard adaptive
change-of-measure and testing inequalities
\citep{bubeck2012survey,lattimore2020bandit}.

\section{Model and Performance Criterion}
\label{sec:model}

\paragraph{Actions and observations.}
Let $(\X,D)$ be a known nonempty compact metric space of diameter at
most one. The horizon $T\ge2$ is an integer. At round $t$, a policy
chooses $X_t\in\X$ using the past and its internal randomization, then
observes both coordinates of a reward vector $(Y_t^1,Y_t^2)$.
For $i\in\{1,2\}$, the conditional mean is $\mu^i(X_t)$ and the
centered noise is conditionally $1$-sub-Gaussian:
\[
 \E\!\left[\exp\{\theta(Y_t^i-\mu^i(X_t))\}
       \mid\mathcal H_{t-1},X_t\right]\le e^{\theta^2/2},
 \qquad \theta\in\R.
\]
The upper bounds need no independence between coordinates.
All lower bounds use independent unit-variance Gaussian coordinates.
The mean functions take values in $[0,1]$ and are $1$-Lipschitz:
\[
 |\mu^i(x)-\mu^i(y)|\le D(x,y),\qquad i\in\{1,2\}.
\]

\paragraph{Lexicographic optimum and regret.}
Continuity and compactness ensure that
\[
 \mu^1_\star=\max_{x\in\X}\mu^1(x),\qquad
 \X_1^\star=\{x:\mu^1(x)=\mu^1_\star\}
\]
are well defined and that a lexicographic optimum exists:
\[
 x^\star\in\argmax_{x\in\X_1^\star}\mu^2(x).
\]
All such optima have the same reward vector. Define the coordinate gaps
and general coordinate regrets by
\[
 \Delta_i(x)=\mu^i(x^\star)-\mu^i(x),\qquad
 R_i(T)=\sum_{t=1}^T\Delta_i(X_t).
\]
The primary gap is nonnegative, while the secondary gap can be negative
on a primary-suboptimal arm. Our upper bounds control the stronger
positive-part quantity
\[
 R_i^+(T)=\sum_{t=1}^T[\Delta_i(X_t)]_+,
 \qquad R_i(T)\le R_i^+(T)\le T.
\]
Thus an upper bound does not rely on cancellation between positive and
negative secondary gaps.

\paragraph{Trade-off information.}
For two objectives, the local trade-off is
\begin{equation}
\label{eq:lambda-def}
 \lambda_\star=
 \inf\left\{\lambda\ge0:
 \mu^2(x)-\mu^2(x^\star)\le\lambda\Delta_1(x),\quad x\in\X
 \right\},
\end{equation}
where the infimum of the empty set is $+\infty$.
It is finite on every finite arm set, but need not have a common
finite upper bound over the instance class. It can also be infinite
on a compact infinite arm space. An envelope is a supplied number
$L\ge0$ satisfying $\lambda_\star\le L$.

\paragraph{Simultaneous expected minimax criterion.}
For an instance class $\F$, let
\begin{equation}
\label{eq:joint-risk}
 \mathfrak M_T(\F)=
 \inf_\pi\sup_{\nu\in\F}
 \max_{i\in\{1,2\}}\E_\nu R_i^\pi(T).
\end{equation}
The same policy must work for every instance and both coordinates.
Policies may use $T$ and the metric, together with any side information
explicitly supplied in the theorem. They do not know the mean rewards.
The expected criterion in \eqref{eq:joint-risk} is distinct from the
high-probability positive-regret criterion defined in
Section~\ref{sec:frontiers}.

For a fixed metric space, $\F(\X)$ denotes all instances satisfying
the observation and Lipschitz assumptions, and we abbreviate
$\mathfrak M_T(\X)=\mathfrak M_T(\F(\X))$. The no-envelope metric
results remain valid if the class is restricted to finite
$\lambda_\star$: the upper bounds hold on the larger class, and all
lower constructions have finite trade-off.
For $K$ arms with the discrete metric $D(a,b)=\ind\{a\ne b\}$, write
$\F_K$ and $\mathfrak M_{T,K}$. With a supplied envelope, write
$\F_K(L)$ and $\mathfrak M_{T,K}(L)$; the infimum is then over policies
allowed to use $L$. Similarly, $\mathfrak M_T(\X;L)$ uses the class
$\F(\X;L)$ and policies given $L$.

\paragraph{Notation.}
The symbols $c,C$ denote positive universal constants that may change
between displays. A subscript indicates permitted dependence.
The notation $\widetilde O$ and $\widetilde\Theta$ suppresses
logarithmic factors in the displayed parameters and horizon, but not
powers of $K$, $L$, $Z$, or $T$. We use natural logarithms except in
$\log_2$. Metric packing and covering conventions are specified in
Section~\ref{sec:profile}.

\section{The Two-Arm Obstruction and Annular Benchmarks}
\label{sec:two-arm}

We restate the two-arm construction of \citet[Theorem~2]{huyuk2021lexicographic}
in the bounded-mean, general-coordinate-regret model used here. Translating
the primary means and fixing a bounded secondary gap changes only numerical
constants. The formulation also makes its consequence for supplied-envelope
annular benchmarks explicit.

\begin{theorem}[Classical two-arm obstruction, reformulated]
\label{thm:two-arm-lower}
There are universal constants $c,T_0>0$ such that, for every $T\ge T_0$, there exist two admissible two-arm Gaussian instances $\nu_0,\nu_1$ for which every policy satisfies
\[
  \max\left\{\E_0R_2(T),\E_1R_1(T)\right\}
  \ge cT^{2/3}.
\]
The null has $\lambda_\star(\nu_0)=0$ and the alternative has a finite value $\lambda_\star(\nu_1)=\Theta(T^{1/3})$.
\end{theorem}

\begin{proof}
Let the arms be $a,b$, fix $g=1/4$, and set $\eps=c_0T^{-1/3}$ for a sufficiently small numerical constant.  On the null,
\[
  \mu_0^1(a)=\mu_0^1(b)=1/2,
  \qquad
  \mu_0^2(a)=g,\quad \mu_0^2(b)=0.
\]
Thus $a$ is lexicographically optimal and $\lambda_\star(\nu_0)=0$.  On the alternative, keep the second coordinate unchanged and set
\[
  \mu_1^1(a)=1/2,
  \qquad
  \mu_1^1(b)=1/2+\eps.
\]
Then $b$ is the unique primary optimum and $\lambda_\star(\nu_1)=g/\eps$.

Let $N_b(T)$ be the number of pulls of $b$.  The relevant regrets are exactly
\[
  R_{2,0}(T)=gN_b(T),
  \qquad
  R_{1,1}(T)=\eps\bigl(T-N_b(T)\bigr).
\]
If $\E_0R_2(T)\ge c\eps T$, we are done.  Otherwise
\[
  \E_0N_b(T)\le \frac{c\eps T}{g}.
\]
For independent unit-variance Gaussian feedback, the adaptive chain rule gives
\[
  \KL(P_0\Vert P_1)
  =\frac{\eps^2}{2}\E_0N_b(T)
  \le \frac{cT\eps^3}{2g}.
\]
Choose $c_0$ and then $c$ so that the last expression is at most $1/32$.  Markov's inequality gives $P_0(N_b(T)>T/2)\le 2c\eps/g$, and Pinsker's inequality transfers the event $N_b(T)\le T/2$ to $P_1$ with constant probability.  On this event, $R_{1,1}(T)\ge\eps T/2$.  Hence
\[
  \max\{\E_0R_2(T),\E_1R_1(T)\}
  \ge c'\eps T
  =\Omega(T^{2/3}).
\]
\end{proof}

The construction also directly separates the known-trade-off zooming benchmark.

\begin{proposition}[Trivial annular data on the hard pair]
\label[proposition]{prop:two-arm-zooming}
For each objective of each instance in \cref{thm:two-arm-lower}, every positive-gap annulus in the sense of \citet{xue2025jmlr} is empty or a singleton.  Hence the annular zooming dimension can be taken as $d_z^i=0$ with coefficient $Z_i=1$.  For the null-secondary pair one may use any constant valid envelope, and for the alternative-primary pair $\Lambda_1(\lambda)\equiv1$.  The corresponding known-trade-off coordinate benchmarks are therefore $\widetilde O(\sqrt T)$.
\end{proposition}

\begin{proof}
The null has no positive primary gaps and only $b$ has a positive
secondary gap. The alternative has only one positive primary gap,
at $a$, and no positive secondary gaps. Every positive-gap annulus
is therefore empty or a singleton. The annular coefficients can be
taken to be one. On the null use the valid envelope $L=1$; on the
alternative use any valid envelope and note that $\Lambda_1(L)=1$.
The supplied-envelope ADLO bound in Table~\ref{tab:related} then
gives $\widetilde O(\sqrt T)$ for the null's secondary coordinate
and for the alternative's primary coordinate.
\end{proof}

This comparison uses two different instance-coordinate pairs. It does
not assert a simultaneous $\widetilde O(\sqrt T)$ supplied-envelope
bound for both coordinates of the large-trade-off alternative.


\begin{corollary}[Polynomial price of adaptation]
\label[corollary]{cor:no-price-false}
No policy independent of $\lambda_\star$ can simultaneously guarantee, on every two-objective instance and for both coordinates, the known-trade-off objective-wise annular-zooming rate up to polylogarithmic factors.
\end{corollary}

\begin{remark}[The lower bound is not a small-$\lambda_\star$ artifact]
\label[remark]{rem:lambda-ge-one}
A three-arm modification makes every hard instance satisfy $\lambda_\star\ge1$.  Let
\[
  (\mu^1(a),\mu^2(a))=(g,g),\quad
  (\mu^1(b),\mu^2(b))=(g,0),\quad
  (\mu^1(c),\mu^2(c))=(0,2g)
\]
on the null, and increase only $\mu^1(b)$ by $\eps$ on the alternative.  The null has $\lambda_\star=1$ and the alternative has $\lambda_\star=g/\eps$.  Moreover,
\[
  R_{1,0}(T)=gN_c(T),
  \qquad
  R_{2,0}(T)=gN_b(T)-gN_c(T),
\]
so
\[
  \max\{\E_0R_1(T),\E_0R_2(T)\}
  \ge \frac g2\E_0N_b(T).
\]
The same information argument yields $\Omega(T^{2/3})$.
\end{remark}

\section{Finite-Arm Minimax Bounds}
\label{sec:finite}

We first give an explicit trade-off-free upper bound.

\begin{algorithm}[t]
\caption{$\LexETC$: finite-arm lexicographic explore-filter-commit}
\label{alg:lex-etc}
\begin{algorithmic}[1]
\Require horizon $T$, arm set $[K]$
\State $\ell\gets\log(4KT)$ and $A\gets K^{1/3}T^{2/3}\ell^{1/3}$
\If{$8A>T$}
  \State play an arbitrary arm for all $T$ rounds and stop
\EndIf
\State $n\gets\left\lceil(T^2\ell/K^2)^{1/3}\right\rceil$, \quad
       $\beta\gets\sqrt{2\ell/n}$
\For{each arm $a\in[K]$}
  \State pull $a$ exactly $n$ times and compute $\widehat\mu^1(a),\widehat\mu^2(a)$
\EndFor
\State $\widehat M_1\gets\max_a\widehat\mu^1(a)$
\State $\widehat\C\gets\{a:\widehat\mu^1(a)\ge \widehat M_1-2\beta\}$
\State $\widehat a\in\argmax_{a\in\widehat\C}\widehat\mu^2(a)$
\State play $\widehat a$ for all remaining rounds
\end{algorithmic}
\end{algorithm}

\begin{theorem}[Finite-arm upper bound]
\label{thm:finite-upper}
There is a universal constant $C$ such that, for every $K\ge2$ and $T\ge2$, $\LexETC$ uses neither $\lambda_\star$ nor an envelope for it and satisfies
\[
  \sup_{\nu\in\F_K}
  \max_{i\le2}\E_\nu R_i(T)
  \le
  C\min\left\{
    T,
    (K\log(4KT))^{1/3}T^{2/3}
  \right\}.
\]
\end{theorem}

\paragraph{Proof idea.}
On the event that every arm-coordinate empirical mean is within $\beta$ of its expectation, a lexicographic optimum remains in $\widehat\C$, every retained arm has primary gap at most $4\beta$, and the committed arm has secondary gap at most $2\beta$.  Thus the positive-part regret is at most $Kn+4T\beta$; the chosen $n$ balances the two terms.  The full proof is in \cref{app:finite-upper}.

The matching lower bound generalizes the two-arm construction.

\begin{theorem}[Finite-arm lower bound]
\label{thm:finite-lower}
There is a universal constant $c>0$ such that, for every $K\ge2$ and $T\ge2$,
\[
  \mathfrak M_{T,K}
  \ge
  c\min\left\{
    T,
    (K-1)^{1/3}T^{2/3}
  \right\}.
\]
\end{theorem}

\paragraph{Proof idea.}
Use one anchor and $N=K-1$ candidate arms.  On the null, all primary means tie and the anchor has a secondary advantage $g$.  Alternative $j$ raises only candidate $j$ by $\eps$ in the primary coordinate.  If null-secondary regret is small, averaging gives a candidate with small null sampling count; the Gaussian KL to that alternative is $\eps^2\E_0N_j/2$.  The information condition is $T\eps^3\lesssim gN$.  Choosing $\eps\asymp\min\{1,(N/T)^{1/3}\}$ yields the result.  See \cref{app:finite-lower}.

\begin{corollary}[Sharp finite-arm expected minimax rate]
\label[corollary]{cor:finite-matching}
For $K\ge2$,
\[
  \boxed{
  \mathfrak M_{T,K}
  =\wtTheta\!\left(
    \min\{T,K^{1/3}T^{2/3}\}
  \right).}
\]
For $K=1$, the value is zero.
\end{corollary}

\section{Supplied Envelopes and the Regret Transition}
\label{sec:envelope}

Suppose the learner is supplied with $L\ge\lambda_\star$, and write
\[
  B=1\vee L.
\]
The key scalarization is
\[
  s_B(x)=2B\mu^1(x)+\mu^2(x),
  \qquad
  G_B(x)=s_B(x^\star)-s_B(x).
\]

\begin{lemma}[Robust-score sandwich]
\label[lemma]{lem:score-sandwich-main}
For every arm $x$,
\[
  G_B(x)\ge B\Delta_1(x),
  \qquad
  \Delta_2(x)\le G_B(x).
\]
In particular, $x^\star$ maximizes $s_B$.
\end{lemma}

\begin{proof}
By \eqref{eq:lambda-def}, $\Delta_2(x)\ge-L\Delta_1(x)$.  Hence
\[
  G_B(x)=2B\Delta_1(x)+\Delta_2(x)
  \ge(2B-L)\Delta_1(x)
  \ge B\Delta_1(x).
\]
The second inequality follows from $G_B(x)-\Delta_2(x)=2B\Delta_1(x)\ge0$.
\end{proof}

A multiscale elimination algorithm estimates both coordinates to accuracy $r_s=2^{-s}$, filters by the empirical robust score, and stops when the prospective regret-weighted exploration cost exceeds $Tr_s$.  It finally selects the surviving arm with the largest empirical secondary mean.  The last step is essential: it makes the commitment secondary gap $O(r_s)$ rather than $O(Br_s)$.

\begin{theorem}[Finite-arm supplied-envelope phase diagram]
\label{thm:finite-envelope}
There are universal constants $c,C>0$ such that, for all $K\ge2$, $T\ge2$, and $L\ge0$,
\[
  c\min\left\{
    T,\sqrt{(1\vee L)(K-1)T},(K-1)^{1/3}T^{2/3}
  \right\}
  \le \mathfrak M_{T,K}(L)
\]
and
\[
  \mathfrak M_{T,K}(L)
  \le
  \wtO\!\left(
  \min\left\{
    T,\sqrt{(1\vee L)KT},K^{1/3}T^{2/3}
  \right\}
  \right).
\]
Hence the bounds match up to logarithmic factors.
\end{theorem}

The proof appears in \cref{app:envelope}. For bounded $L$, a
primary-only construction gives the lower bound. For larger $L$, the
same hidden-candidate family uses $g=\min\{1/4,L\eps\}$ and covers
both the square-root and cubic branches. Table~\ref{tab:regimes}
makes the transition explicit. The thresholds below suppress
logarithmic factors; the theorem itself includes the linear cap for
every $K$ and $T$.

\begin{table}[t]
\centering
\small
\renewcommand{\arraystretch}{1.15}
\caption{Finite-arm regimes for the joint expected criterion.
Here $B=1\vee L$. Rates are up to logarithms and universal constants.}
\label{tab:regimes}
\begin{tabularx}{\textwidth}{@{}>{\raggedright\arraybackslash}p{5.1cm}>{\raggedright\arraybackslash}p{3.4cm}Y@{}}
\toprule
Parameter range & Minimax rate & Source of difficulty \\
\midrule
$K\ge T$ & $T$ & Too many potential optima for the horizon. \\
$2\le K\le T$, $1\le B\le(T/K)^{1/3}$
& $\sqrt{BKT}$ & The envelope bounds secondary search losses. \\
$2\le K\le T$, $B\ge(T/K)^{1/3}$
& $K^{1/3}T^{2/3}$ & The fixed cost of searching primary ties dominates. \\
No supplied envelope, $K\ge2$
& $\min\{T,K^{1/3}T^{2/3}\}$ & The policy must cover arbitrarily large finite trade-offs. \\
\bottomrule
\end{tabularx}
\end{table}

\section{Primary-Tie Geometry and Adaptation}
\label{sec:profile}

The finite-arm lower bound counts possible hidden primary improvements.  In a metric space, the corresponding quantity is the packing entropy of arms statistically compatible with a primary tie.

\begin{definition}[Primary near-optimality packing profile]
\label[definition]{def:profile}
Fix $c_{\mathrm{nopt}}=64$.  For $0<r\le1$, define
\[
  \Pcal_\nu(r)
  =\Pack\!\left(
    \{x\in\X:\Delta_1(x)\le c_{\mathrm{nopt}}r\},
    r
  \right).
\]
For $d\ge0$ and $Z\ge1$, say that $\nu$ has power profile $(d,Z)$ if
\[
  \Pcal_\nu(r)\le Zr^{-d},\qquad 0<r\le1.
\]
\end{definition}

This is the classical near-optimality packing profile from $X$-armed bandits.  The contribution is to identify it as the missing complexity for unknown lexicographic trade-offs.

For a subset $U\subseteq\X$, let $\Pack(U,r)$ be the largest
cardinality of a subset whose distinct points are at distance at least
$r$, and let $\Cov(U,r)$ be the minimum number of closed radius-$r$
balls needed to cover $U$, with both numbers set to zero for the
empty set. Compactness of $\X$ makes both numbers finite.
We assume a measurable maximal-net oracle that, given a nonempty compact
$U$ and $r>0$,
returns a finite inclusion-maximal $r$-separated subset. Such a subset
covers $U$ by closed radius-$r$ balls. The oracle's computational cost
is not included in the observation budget. A supplied covering
hierarchy is discussed in \cref{app:entropy}.

\begin{algorithm}[t]
\caption{$\PFETC$: profile-oblivious primary filtering}
\label{alg:pf-etc}
\begin{algorithmic}[1]
\Require horizon $T$, confidence $\delta$, maximal-net oracle
\State $U_{-1}\gets\X$, $\Gamma_{-1}\gets0$, $r_s\gets2^{-s}$
\For{$s=0,1,2,\ldots$}
  \State $A_s\gets$ a measurable maximal $r_s$-separated subset of $U_{s-1}$
  \State $n_s\gets\left\lceil2r_s^{-2}\log\!\left(\frac{8\pi^2(s+1)^2\max\{2,|A_s|\}}{3\delta}\right)\right\rceil$
  \State $\Gamma_s^{\rm pro}\gets\Gamma_{s-1}+64|A_s|n_s$
  \If{$\Gamma_s^{\rm pro}>Tr_s$}
    \State stop before sampling stage $s$
  \EndIf
  \State sample every center in $A_s$ exactly $n_s$ fresh times
  \State $G_s\gets\{x\in A_s:\widehat\mu_s^1(x)\ge\max_{y\in A_s}\widehat\mu_s^1(y)-3r_s\}$
  \State $U_s\gets U_{s-1}\cap\bigcup_{x\in G_s}\overline B(x,r_s)$ and $\Gamma_s\gets\Gamma_s^{\rm pro}$
\EndFor
\If{no stage was completed}
  \State play a fixed arm for all remaining rounds
\Else
  \State let $\tau$ be the last completed stage and commit to
  $\widehat x_\tau\in\argmax_{x\in G_\tau}\widehat\mu_\tau^2(x)$
\EndIf
\end{algorithmic}
\end{algorithm}

\begin{theorem}[Profile-oblivious oracle inequality]
\label{thm:profile-oracle}
Algorithm \ref{alg:pf-etc} uses neither $\lambda_\star$ nor any profile parameter.  For
\[
  \overline n_s(p,\delta)
  =\left\lceil2r_s^{-2}\log\!\left(
  \frac{8\pi^2(s+1)^2\max\{2,p\}}{3\delta}
  \right)\right\rceil
\]
and
\[
  H_U(s;\nu,\delta)
  =64\sum_{j=0}^s
  \Pcal_\nu(r_j)\overline n_j(\Pcal_\nu(r_j),\delta),
\]
there is a universal constant $C$ such that
\[
  \max_{i\le2}\E_\nu R_i(T)
  \le
  C\left[
    \min\left\{T,\inf_{s\ge0}
    \bigl(H_U(s;\nu,\delta)+Tr_s\bigr)\right\}
    +\delta T
  \right].
\]
\end{theorem}

\begin{corollary}[Power-profile upper bound]
\label[corollary]{cor:power-upper}
If $\Pcal_\nu(r)\le Zr^{-d}$ and $\delta=T^{-2}$, then the same policy, without knowing $d$ or $Z$, satisfies
\[
  \max_{i\le2}\E_\nu R_i(T)
  \le
  \wtO_d\!\left(
  \min\left\{
    T,
    Z^{1/(d+3)}T^{(d+2)/(d+3)}
  \right\}
  \right).
\]
\end{corollary}

The proof is in \cref{app:profile-upper}.  On the high-probability event, every active net contains a representative of $x^\star$, retained centers have primary gap $O(r_s)$, and the final empirical-secondary maximizer has both coordinate gaps $O(r_s)$.  The stopping rule compares the realized exploration cost to $Tr_s$ and thereby adapts to the unknown profile.

\subsection{Matching regular metric classes}

To state a matching lower bound, one needs a sufficiently rich metric class; a profile upper bound alone cannot force hardness on every subclass.

\begin{assumption}[Regular anchor-plus-tie space]
\label[assumption]{ass:regular}
The arm space is $\X=\{a\}\cup S$, with $D(a,z)=1$ for $z\in S$
and $\operatorname{diam}(S)\le1/4$. Fix $d\ge0$, $c_->0$, $c_+>0$,
and $r_0\in(0,1/16]$. For some $Z\ge1$,
\[
 \Pack(S,4r)\ge c_-Zr^{-d},\qquad
 \Cov(S,r)\le c_+Zr^{-d},\qquad 0<r\le r_0.
\]
\end{assumption}

\begin{theorem}[Sharp regular-metric adaptation rate]
\label{thm:regular-metric}
For every fixed metric satisfying \cref{ass:regular},
\[
 \mathfrak M_T(\X)
 =\widetilde\Theta_{d,c_-,c_+,r_0}\!\left(
 \min\{T,Z^{1/(d+3)}T^{(d+2)/(d+3)}\}\right).
\]
The constants are uniform in $T$ and $Z$ when the displayed regularity
parameters are fixed. The conclusion also holds on the subclass of
finite local trade-offs. The upper policy needs neither
$\lambda_\star$ nor the regularity parameters.
\end{theorem}

The lower bound chooses $N\asymp Z\eps^{-d}$ disjoint hidden locations.  The null makes all of $S$ primary-tied and assigns a secondary advantage to the anchor.  Alternative $v$ inserts the cone $(\eps-D(\cdot,v))_+$ in the primary objective.  The information balance becomes
\[
  T\eps^{d+3}\lesssim Z,
\]
which yields the $d+3$ denominator.  Complete details are in \cref{app:metric-lower}.

\subsection{Arbitrary tie-metric entropy}
\label{sec:entropy}

The power-law assumption is not essential.  Fix $\X=\{a\}\cup S$ with $D(a,S)=1$ and arbitrary nonempty compact $S$ of diameter at most $1/4$.  Let
\[
  r_s=2^{-s-4},
  \qquad
  M_s=\Pack(S,r_s),
  \qquad
  s_T=\left\lceil\log_2(16T)\right\rceil+2.
\]
Define
\[
  \Psi_T^\infty(S)
  =T\max_{0\le s\le s_T}
  \min\left\{r_s,\left(\frac{M_s}{T}\right)^{1/3}\right\},
\]
and, for $B\ge1$,
\[
  \Psi_T^B(S)
  =T\max_{0\le s\le s_T}
  \min\left\{
    r_s,
    \sqrt{\frac{BM_s}{T}},
    \left(\frac{M_s}{T}\right)^{1/3}
  \right\}.
\]

\begin{theorem}[Arbitrary tie-metric entropy]
\label{thm:entropy}
There are universal constants $c,C>0$ such that, for every $T\ge2$,
\[
  c\Psi_T^\infty(S)
  \le \mathfrak M_T(\X)
  \le C\log^{2/3}(eT)\Psi_T^\infty(S).
\]
If an envelope $L$ is supplied and $B=1\vee L$, then
\[
  c\Psi_T^B(S)
  \le \mathfrak M_T(\X;L)
  \le C\log(eT)\Psi_T^B(S).
\]
The upper policy does not know the entropy sequence $(M_s)_s$.
\end{theorem}

This result is metric-wise for the anchor-plus-tie geometry.  It is not a claim of polynomial-time net construction from an arbitrary black-box metric.  If a well-shaped hierarchical covering tree is supplied, the traversal is output-sensitive and linear in the number of generated nodes.  The proof is in \cref{app:entropy}.

\section{Expected-Regret Trade-offs and Confidence Bounds}
\label{sec:frontiers}

The balanced minimax lower bound follows from a more detailed constraint on how a policy distributes regret between the null and the alternatives.

Consider one anchor $a$ and $N$ candidates.  Fix $g\in(0,1/4]$.  On the null, all primary means tie, the anchor has secondary mean $g$, and each candidate has secondary mean zero.  Alternative $j$ raises only candidate $j$ by $\eps$ in the primary coordinate.  Define
\[
  r_0=\E_0R_2(T),
  \qquad
  r_j(\eps)=\E_j^\eps R_1(T),
  \qquad
  r_\star(\eps)=\max_{j\in[N]}r_j(\eps).
\]

\begin{theorem}[Expected-regret Pareto lower constraint]
\label{thm:pareto}
For every policy and every $0<\eps\le g$,
\[
  \boxed{
  \frac{r_0}{gN}
  +\frac{r_\star(\eps)}{\eps}
  \ge
  \frac T2
  \exp\left(-\frac{\eps^2r_0}{2gN}\right).}
\]
Consequently, if $r_0\le gT/8$,
\[
  \sup_{0<\eps\le g/2}r_\star(\eps)
  \ge
  cT\min\left\{
    g,
    \sqrt{\frac{gN}{r_0+gN}}
  \right\}.
\]
In the square-root branch,
\[
  (r_0+gN)
  \left(\sup_{\eps}r_\star(\eps)\right)^2
  \ge cgNT^2.
\]
\end{theorem}

This theorem is a Pareto \emph{lower constraint}; we do not claim a family of algorithms matching every point of the boundary.  The proof is in \cref{app:pareto}.

For high-probability guarantees, define
\[
  \mathfrak Q_{T,K}(\delta)
  =\inf_\pi\inf\left\{r:
  \sup_{\nu\in\F_K}
  \Pp_\nu\left(\max_{i\le2}R_i^+(T)>r\right)
  \le\delta
  \right\}.
\]

\begin{theorem}[General finite-arm high-probability lower bound]
\label{thm:hp-lower}
There are universal constants $c,C>0$ such that, for $K=N+1\ge2$, $T\ge C$, and $0<\delta\le1/16$,
\[
  \mathfrak Q_{T,K}(\delta)
  \ge
  c\min\left\{
    T,
    \bigl(N\log(1/\delta)\bigr)^{1/3}T^{2/3}
  \right\}.
\]
\end{theorem}

\begin{theorem}[Sharp two-arm high-probability rate]
\label{thm:two-hp}
There are universal constants $c,C>0$ such that, for $T\ge C$ and $0<\delta\le1/16$,
\[
  c\min\left\{
    T,T^{2/3}\log^{1/3}(1/\delta)
  \right\}
  \le \mathfrak Q_{T,2}(\delta)
\]
\[
  \le
  C\min\left\{
    T,T^{2/3}\log^{1/3}(8/\delta)
  \right\}.
\]
\end{theorem}

\begin{proposition}[General-$K$ high-probability bracket]
\label[proposition]{prop:k-hp}
For every $K\ge2$, $T\ge C$, and $0<\delta\le1/16$,
\[
  c\min\left\{
    T,
    ((K-1)\log(1/\delta))^{1/3}T^{2/3}
  \right\}
  \le \mathfrak Q_{T,K}(\delta)
\]
and
\[
  \mathfrak Q_{T,K}(\delta)
  \le
  C\min\left\{
    T,
    (K\log(4K/\delta))^{1/3}T^{2/3}
  \right\}.
\]
Thus the confidence dependence is sharp; the current general-$K$ upper bound has an additional $\log^{1/3}K$ factor.
\end{proposition}

The high-probability lower bound uses a stopped-history KL argument: either the null explores many candidates with probability at least $\delta$, or some candidate is under-sampled with constant null probability; stopping at its $m$-th pull keeps KL at $O(\log(1/\delta))$ and binary-KL inversion transfers the under-sampling event.  Proofs are in \cref{app:high-prob}.

\section{Discussion}
\label{sec:discussion}

The results distinguish three sources of statistical difficulty. Ordinary
primary-reward uncertainty gives the square-root lower bound. A bounded
trade-off permits a larger secondary loss near a primary optimum and
increases that rate by a square-root dependence on the envelope. When
the envelope is sufficiently large, the fixed loss of searching a primary
tie dominates; further increases in the envelope do not change the
finite-horizon minimax rate.

The metric results make the search region explicit. The profile upper
bound depends on the primary near-optimal set of the particular instance.
Its matching lower bounds concern rich classes, constructed by allowing
a primary improvement at any of many separated tie locations. A small
profile therefore supplies an upper bound; it does not, by itself, imply
that every instance with that profile is hard. The arbitrary-entropy
theorem is likewise a minimax statement on a fixed anchor-plus-tie
space, rather than an instance-wise equivalence.

Several questions remain. The general-$K$ high-probability bounds leave
a logarithmic gap in $K$. The Pareto inequality constrains what a policy
can achieve, but a family of algorithms attaining every point of that
constraint is not established here. More than two objectives may require
several interacting search regions. Finally, the metric algorithms assume
measurable maximal-net access; their oracle guarantees do not supply an
efficient representation of an arbitrary compact metric space.

\section{Conclusion}

Primary ties link the classical impossibility of fast prior-free
lexicographic learning to a quantitative theory of simultaneous regret.
For finite arms, the sharp dependence is $K^{1/3}T^{2/3}$ without an
envelope, and $\min\{\sqrt{(1\vee L)KT},K^{1/3}T^{2/3}\}$ with a
supplied envelope, in both cases capped by $T$ and up to logarithms.
For metric arms, the packing profile of primary near-optimal actions
governs a profile-adaptive upper bound and yields matching $d+3$ rates
on regular tie spaces. The entropy and confidence refinements describe
the same testing cost at finer resolution.

\appendix

\section{Finite-Arm Proofs}
\label[appendix]{app:finite-upper}

\begin{lemma}[Terminal accuracy of $\LexETC$]
\label[lemma]{lem:finite-terminal}
On the event
\[
  \mathcal E=
  \left\{
  |\widehat\mu^i(a)-\mu^i(a)|\le\beta,
  \ \forall a\in[K],\ i\in\{1,2\}
  \right\},
\]
the arm committed by Algorithm~\ref{alg:lex-etc} satisfies
\[
  \Delta_1(\widehat a)\le4\beta,
  \qquad
  \Delta_2(\widehat a)\le2\beta.
\]
\end{lemma}

\begin{proof}
Let $a^\star$ be a lexicographic optimum.  On $\mathcal E$,
\[
  \widehat M_1-\widehat\mu^1(a^\star)
  \le(\mu^1_\star+\beta)-(\mu^1_\star-\beta)
  =2\beta,
\]
so $a^\star\in\widehat\C$.  For every $a\in\widehat\C$,
\[
  \widehat\mu^1(a)
  \ge \widehat\mu^1(a^\star)-2\beta
  \ge \mu^1_\star-3\beta,
\]
and hence $\mu^1(a)\ge\mu^1_\star-4\beta$.  Since $\widehat a$ maximizes the empirical secondary mean over a set containing $a^\star$,
\[
  \mu^2(\widehat a)
  \ge\widehat\mu^2(\widehat a)-\beta
  \ge\widehat\mu^2(a^\star)-\beta
  \ge\mu^2(a^\star)-2\beta.
\]
\end{proof}

\begin{proof}[Proof of \cref{thm:finite-upper}]
By sub-Gaussian concentration,
\[
  \Pp(\mathcal E^c)
  \le4K\exp(-n\beta^2/2)
  =4Ke^{-\ell}
  \le T^{-1}.
\]
Every positive instantaneous coordinate gap is at most one.  Thus exploration contributes at most $Kn$ to either positive-part regret.  On $\mathcal E$, \cref{lem:finite-terminal} bounds commitment by $4T\beta$; on $\mathcal E^c$, regret is at most $T$.  Therefore
\[
  \max_i\E R_i(T)
  \le Kn+4T\beta+1.
\]
In the nontrivial branch, let $A=K^{1/3}T^{2/3}\ell^{1/3}\le T/8$.  Cubing gives $K\ell\le T/512$, hence $K\le T/512$.  Since
\[
  n\le(T^2\ell/K^2)^{1/3}+1,
\]
we have $Kn\le A+K<T/2$, so the algorithm is feasible.  Moreover,
\[
  Kn\le C A,
  \qquad
  T\beta\le C A.
\]
In the trivial branch, regret is at most $T$, which is within a universal constant of the linearly capped target.
\end{proof}

\section{Finite-Arm Lower Bounds}
\label[appendix]{app:finite-lower}

Let the arms be $\{a,1,\ldots,N\}$, $N=K-1$, and let $0<g\le1/4$.  On the null,
\[
  \mu_0^1(x)=1/2\quad\forall x,
  \qquad
  \mu_0^2(a)=g,
  \qquad
  \mu_0^2(j)=0.
\]
For alternative $j$, keep the second coordinate unchanged and set
\[
  \mu_j^1(j)=1/2+\eps,
  \qquad
  \mu_j^1(x)=1/2\quad(x\ne j).
\]
Use independent unit-variance Gaussian noise.

\begin{lemma}[Costly search versus a hidden candidate]
\label[lemma]{lem:finite-costly}
There is a universal $c_0>0$ such that, if
\[
  0<\eps\le g/2,
  \qquad
  T\eps^3\le c_0gN,
\]
then every policy satisfies
\[
  \max\left\{
    \E_0R_2(T),
    \max_{j\in[N]}\E_jR_1(T)
  \right\}
  \ge c_0\eps T.
\]
\end{lemma}

\begin{proof}
Let $N_S(T)=\sum_jN_j(T)$.  Under the null,
\[
  R_2(T)=gN_S(T).
\]
If $\E_0R_2(T)\ge c_0\eps T$, we are done.  Otherwise, some $j$ satisfies
\[
  \E_0N_j(T)
  \le\frac{c_0\eps T}{gN}.
\]
The adaptive Gaussian chain rule gives
\[
  \KL(P_0\Vert P_j)
  =\frac{\eps^2}{2}\E_0N_j(T)
  \le\frac{c_0T\eps^3}{2gN}.
\]
Choose constants so that this is at most $1/128$.  Let
$E_j=\{N_j(T)\le T/2\}$.  Markov's inequality gives
\[
  P_0(E_j^c)
  \le\frac{2\E_0N_j(T)}{T}
  \le\frac{2c_0\eps}{gN}
  \le\frac18.
\]
Pinsker's inequality therefore implies $P_j(E_j)\ge1/2$.  On $E_j$, at least $T/2$ pulls have primary gap $\eps$, so
\[
  \E_jR_1(T)\ge\frac{\eps T}{4}.
\]
\end{proof}

\begin{proof}[Proof of \cref{thm:finite-lower}]
Fix $g=1/4$ and take
\[
  \eps=c\min\left\{1,\left(\frac NT\right)^{1/3}\right\}
\]
with $c$ sufficiently small.  Then the conditions of \cref{lem:finite-costly} hold and
\[
  \mathfrak M_{T,K}
  \ge c'\eps T
  \ge c''\min\{T,N^{1/3}T^{2/3}\}.
\]
\end{proof}

\section{Proof of the Supplied-Envelope Phase Diagram}
\label[appendix]{app:envelope}

We first make the robust-score algorithm precise and then prove the lower
bound with a single construction that covers the entire range $L\ge2$.
Throughout this section, $B=1\vee L$, $r_s=2^{-s}$, and
$q_s=\min\{1,Br_s\}$.

\subsection{Robust-score elimination}

Fix a confidence level $\delta\in(0,1)$.  Set $A_0=[K]$ and
$\Gamma_{-1}=0$.  For a nonempty active set $A_s$, define
\[
 n_s=\left\lceil 2r_s^{-2}
 \log\left(\frac{8\pi^2(s+1)^2\max\{2,|A_s|\}}{3\delta}\right)
 \right\rceil,
 \qquad
 \Gamma_s^{\mathrm{pro}}=
 \Gamma_{s-1}+64q_s|A_s|n_s.
\]
Before taking any stage-$s$ samples, compare
$\Gamma_s^{\mathrm{pro}}$ with $Tr_s$.
If $\Gamma_s^{\mathrm{pro}}>Tr_s$, stop exploration.  When $s=0$, play
an arbitrary arm for the entire horizon.  When $s\ge1$, set $\tau=s-1$
and commit for all remaining rounds to
\[
 \widehat a\in\argmax_{a\in A_{\tau+1}}
                     \widehat\mu^2_\tau(a),
\]
where the empirical means were computed at the last completed stage.
Otherwise, pull each arm in $A_s$ exactly $n_s$ times, using fresh samples,
set $\Gamma_s=\Gamma_s^{\mathrm{pro}}$, and form
\[
 \widehat s_{B,s}(a)=2B\widehat\mu^1_s(a)+\widehat\mu^2_s(a),
 \qquad
 A_{s+1}=\left\{a\in A_s:
 \widehat s_{B,s}(a)\ge
 \max_{b\in A_s}\widehat s_{B,s}(b)-12Br_s\right\}.
\]
The empirical maximizer always survives, so every active set is nonempty.

The algorithm is feasible on every sample path.  Indeed, if stage $s$ is
accepted, its cumulative sample count, including that stage, satisfies
\begin{equation}
\label{eq:envelope-feasibility}
 \sum_{j=0}^s |A_j|n_j
 \le \frac{\Gamma_s}{64q_s}
 \le \frac{Tr_s}{64q_s}
 \le \frac{T}{64}.
\end{equation}
Here $q_j\ge q_s\ge r_s$ for $j\le s$.  Thus there is enough time for
every accepted stage and for the final commitment.  The algorithm also
terminates: $n_s\ge r_s^{-2}$ and $q_s\ge r_s$ imply
$\Gamma_s^{\mathrm{pro}}\ge64/r_s$, so sufficiently fine stages must fail
the stopping test.  In particular, no stage with
$s>\lceil\log_2 T\rceil$ can be reached.

\begin{lemma}[Confidence and robust-score invariants]
\label[lemma]{lem:finite-score-inv}
With probability at least $1-\delta$, all completed-stage empirical means
satisfy
\[
 |\widehat\mu_s^i(a)-\mu^i(a)|\le r_s,
 \qquad a\in A_s,\quad i\in\{1,2\}.
\]
On this event, a lexicographic optimum is never eliminated.  Every arm
active at stage $s\ge1$ satisfies
\[
 \Delta_1(a)\le36r_s,
 \qquad [\Delta_2(a)]_+\le36Br_s.
\]
If at least one stage is completed, the committed arm satisfies
\[
 \Delta_1(\widehat a)\le18r_\tau,
 \qquad \Delta_2(\widehat a)\le2r_\tau.
\]
\end{lemma}

\begin{proof}
Conditional on the history before stage $s$, the active set is fixed and
the fresh observations satisfy the assumed conditional sub-Gaussian
bound.  A union bound over its two coordinates and all its arms gives a
conditional failure probability at most
\[
 4|A_s|\exp(-n_sr_s^2/2)
 \le\frac{3\delta}{2\pi^2(s+1)^2}.
\]
Summing over stages gives a total failure probability at most $\delta$.
This argument also applies when a stage is not reached, by assigning it
zero failure probability.

On the confidence event, the score error is at most
$(2B+1)r_s\le3Br_s$.  By \cref{lem:score-sandwich-main}, a lexicographic
optimum $a^\star$ maximizes the true score.  Its empirical score is within
$6Br_s$ of the empirical maximum, so it survives the $12Br_s$ filter.
Every survivor $a\in A_{s+1}$ satisfies
\[
 G_B(a)
 \le 3Br_s+12Br_s+3Br_s=18Br_s.
\]
The robust-score sandwich gives $\Delta_1(a)\le18r_s$ and
$[\Delta_2(a)]_+\le18Br_s$.  Replacing $r_{s-1}$ by $2r_s$ proves the
bounds for arms active at the beginning of stage $s$.
Finally, $a^\star\in A_{\tau+1}$, so empirical secondary maximization
at stage $\tau$ gives
\[
 \mu^2(\widehat a)
 \ge\widehat\mu^2_\tau(\widehat a)-r_\tau
 \ge\widehat\mu^2_\tau(a^\star)-r_\tau
 \ge\mu^2(a^\star)-2r_\tau.
\]
Its primary gap is bounded by the survivor estimate from stage $\tau$.
\end{proof}

\begin{lemma}[Finite envelope upper bound]
\label[lemma]{lem:finite-envelope-upper}
There is a universal constant $C$ such that, with
$\ell=\log(64KT/\delta)$, robust-score elimination satisfies, with
probability at least $1-\delta$,
\[
 \max_{i\le2}R_i^+(T)
 \le C\min\left\{
 T,\sqrt{BK\ell T},(K\ell)^{1/3}T^{2/3}
 \right\}.
\]
In particular, choosing $\delta=1/T$ proves the upper bound in
\cref{thm:finite-envelope}.
\end{lemma}

\begin{proof}
Work on the confidence event in \cref{lem:finite-score-inv}.  A pull in
stage $s\ge1$ has either positive coordinate gap at most
$\min\{1,36Br_s\}\le36q_s$.  The same bound holds at stage zero because
$q_0=1$.  Thus the exploration regret is at most
$(36/64)\Gamma_\tau$.  If the first failed stage is $j\ge1$, then
$\tau=j-1$, and the commitment guarantee gives
\begin{equation}
\label{eq:envelope-crossing-risk}
 \max_{i\le2}R_i^+(T)
 \le\Gamma_\tau+18Tr_\tau
 \le19Tr_\tau=38Tr_j.
\end{equation}
If $j=0$, the deterministic bound $R_i^+(T)\le T$ applies.

We next relate this stopping rule to a deterministic cost bound.
All attempted stages satisfy $s\le\lceil\log_2 T\rceil$.  Since
$|A_s|\le K$, their sample sizes obey
$n_s\le C_0r_s^{-2}\ell$ for a universal $C_0$.  Define, for every
integer $s\ge0$,
\[
 H_s=64C_0K\ell\sum_{u=0}^s q_ur_u^{-2}.
\]
For each attempted stage, $\Gamma_s^{\mathrm{pro}}\le H_s$.
The sequence $H_s$ is nondecreasing.  At the first failed stage $j$,
\[
 H_j\ge\Gamma_j^{\mathrm{pro}}>Tr_j.
\]
For $s\le j$, we have $Tr_s\ge Tr_j$; for $s\ge j$, we have
$H_s\ge H_j>Tr_j$.  Hence
\[
 \inf_{s\ge0}\{H_s+Tr_s\}\ge Tr_j.
\]
Combined with \eqref{eq:envelope-crossing-risk}, this gives an upper
bound of $38\inf_s\{H_s+Tr_s\}$ when $j\ge1$.  When $j=0$,
$H_s\ge H_0>T$ for all $s$, and the same conclusion follows from the
trivial bound $T$.

Geometric summation yields
\[
 \sum_{u=0}^s q_ur_u^{-2}
 \le\min\left\{\sum_{u=0}^s r_u^{-2},
                  B\sum_{u=0}^s r_u^{-1}\right\}
 \le2\min\{r_s^{-2},Br_s^{-1}\}.
\]
Consequently, the risk is at most a universal constant times
\[
 \inf_{s\ge0}\left\{
 K\ell\min\{r_s^{-2},Br_s^{-1}\}+Tr_s\right\}.
\]
Choosing dyadic scales within a factor two of
$(K\ell/T)^{1/3}$ and $(BK\ell/T)^{1/2}$, whenever these are at most
one, proves the two nonlinear branches.  The deterministic bound $T$
supplies the remaining branch.  Finally, the confidence failure adds at
most $T\delta$ in expectation; with $\delta=1/T$, this is absorbed into
the displayed bound.
\end{proof}

\subsection{Lower bounds for all envelope sizes}

We include the primary-only lower bound to make the envelope argument
self-contained.

\begin{lemma}[Primary-only finite-arm lower bound]
\label[lemma]{lem:primary-only-finite}
There is a universal $c>0$ such that, for $K\ge2$ and $T\ge2$, every
policy has expected primary regret at least
$c\min\{T,\sqrt{KT}\}$ on some $K$-arm Gaussian instance with constant
secondary means.  In particular, every such instance has
$\lambda_\star=0$.
\end{lemma}

\begin{proof}
Let $P_0$ denote the experiment with all primary means equal to $1/2$.
For $j\in[K]$, let $P_j$ raise only arm $j$ to $1/2+\eta$, where
\[
 \eta=c_1\min\{1,\sqrt{K/T}\},\qquad 0<c_1\le1/4.
\]
All coordinates have independent unit-variance Gaussian noise and the
secondary means are zero.  The adaptive chain rule and Pinsker's
inequality give
\[
 \KL(P_0\Vert P_j)=\frac{\eta^2}{2}\E_0N_j(T),
 \qquad
 \E_jN_j(T)\le\E_0N_j(T)
                  +\frac{T\eta}{2}\sqrt{\E_0N_j(T)}.
\]
Averaging over $j$ and applying Cauchy--Schwarz yields
\[
 \frac1K\sum_{j=1}^K\E_jN_j(T)
 \le\frac TK+\frac{T\eta}{2}\sqrt{\frac TK}.
\]
Since $K\ge2$ and $\eta\sqrt{T/K}\le c_1$, the average primary
regret on these alternatives is at least
\[
 \frac1K\sum_{j=1}^K\E_jR_1(T)
 =\frac{\eta}{K}\sum_{j=1}^K\bigl(T-\E_jN_j(T)\bigr)
 \ge\frac{\eta T}{4}.
\]
At least one alternative therefore has the claimed regret.
\end{proof}

\begin{lemma}[Envelope lower bound]
\label[lemma]{lem:envelope-lower}
For $N=K-1\ge1$, $T\ge2$, and $L\ge0$,
\[
 \mathfrak M_{T,K}(L)
 \ge c\min\left\{
 T,\sqrt{(1\vee L)NT},N^{1/3}T^{2/3}
 \right\}
\]
for a universal constant $c>0$.
\end{lemma}

\begin{proof}
For $0\le L<2$, \cref{lem:primary-only-finite} gives
$c\min\{T,\sqrt{NT}\}$.  This is within a universal factor of the
claimed minimum: when $N\le T$, its square-root branch is at most
$\sqrt{2NT}$, and when $N>T$ the target is at most $T$.

Now suppose $L\ge2$.  Use the anchor-and-candidates family from
\cref{app:finite-lower}, with
\[
 \eps=c_2\min\left\{1,\sqrt{\frac{LN}{T}},
                         \left(\frac NT\right)^{1/3}\right\},
 \qquad
 g=\min\{1/4,L\eps\}.
\]
Here $c_2>0$ is a sufficiently small universal constant, in particular
$c_2\le1/8$.  The null has $\lambda_\star=0$, and each alternative
has
\[
 \lambda_\star=\frac g\eps\le L.
\]
The inequality $\eps\le g/2$ holds in both cases: if $g=L\eps$ it
follows from $L\ge2$, and if $g=1/4$ it follows from $\eps\le1/8$.
When $g=L\eps$, we have
\[
 \frac{T\eps^3}{gN}
 =\frac{T\eps^2}{LN}\le c_2^2.
\]
When $g=1/4$, we instead have
\[
 \frac{T\eps^3}{gN}\le4c_2^3.
\]
Choosing $c_2$ so that both upper bounds are at most the constant in
\cref{lem:finite-costly}, we may apply that lemma.  The result is a
lower bound $c'\eps T$, which is exactly the claimed three-branch
minimum up to a universal constant.
\end{proof}

\begin{proof}[Proof of \cref{thm:finite-envelope}]
Combine \cref{lem:finite-envelope-upper,lem:envelope-lower} and use
$K/2\le K-1\le K$ for $K\ge2$.
\end{proof}

\section{Proof of the Profile-Oblivious Oracle Inequality}
\label[appendix]{app:profile-upper}

\begin{lemma}[Uniform confidence for random stagewise nets]
\label[lemma]{lem:stage-confidence}
Conditionally on the past, let $A_s$ be a finite random set and pull each center $n_s$ fresh times as in Algorithm~\ref{alg:pf-etc}.  Simultaneously over all sampled stages, centers, and coordinates,
\[
  |\widehat\mu_s^i(x)-\mu^i(x)|\le r_s
\]
holds with probability at least $1-\delta$.
\end{lemma}

\begin{proof}
Conditionally on the history before stage $s$, $A_s$ is fixed.  A union bound over $2|A_s|$ coordinate-center pairs gives conditional failure probability at most
\[
  4|A_s|e^{-n_sr_s^2/2}
  \le\frac{3\delta}{2\pi^2(s+1)^2}.
\]
Summing over $s$ proves the claim.
\end{proof}

\begin{lemma}[Primary-filtering invariants]
\label[lemma]{lem:profile-invariants}
On the event of \cref{lem:stage-confidence}, for every completed stage $s$:
\begin{enumerate}[label=(\roman*)]
  \item $x^\star\in U_s$;
  \item every $x\in G_s$ satisfies $\Delta_1(x)\le6r_s$;
  \item $U_s\subseteq\{x:\Delta_1(x)\le7r_s\}$;
  \item every completed candidate net, and the first rejected candidate net, satisfies $|A_s|\le\Pcal_\nu(r_s)$;
  \item the terminal arm satisfies
  $\Delta_1(\widehat x_\tau)\le6r_\tau$ and
  $\Delta_2(\widehat x_\tau)\le3r_\tau$.
\end{enumerate}
\end{lemma}

\begin{proof}
Assume inductively that $x^\star\in U_{s-1}$.  By maximality, some $\bar x_s\in A_s$ satisfies $D(\bar x_s,x^\star)\le r_s$.  On the confidence event,
\[
  \widehat\mu_s^1(\bar x_s)
  \ge\mu^1(x^\star)-2r_s,
  \qquad
  \max_{x\in A_s}\widehat\mu_s^1(x)
  \le\mu^1(x^\star)+r_s.
\]
Hence $\bar x_s\in G_s$ and $x^\star\in U_s$.  For $x\in G_s$,
\[
  \widehat\mu_s^1(x)
  \ge\widehat\mu_s^1(\bar x_s)-3r_s
  \ge\mu^1(x^\star)-5r_s,
\]
so $\Delta_1(x)\le6r_s$.  Lipschitzness gives the active-set containment.  For $s\ge1$, $A_s\subseteq\{\Delta_1\le14r_s\}$ and is $r_s$-separated, so $|A_s|\le\Pcal_\nu(r_s)$ because $c_{\mathrm{nopt}}=64$.  At $s=0$, the claim is immediate from bounded gaps.  Finally, $\bar x_\tau\in G_\tau$, and empirical secondary maximization against it gives
\[
  \mu^2(\widehat x_\tau)
  \ge\mu^2(\bar x_\tau)-2r_\tau
  \ge\mu^2(x^\star)-3r_\tau.
\]
\end{proof}

\begin{lemma}[Deterministic crossing]
\label[lemma]{lem:crossing}
Let $H_s$ be nondecreasing and $r_s=2^{-s}$.  If
\[
  s_\star=\min\{s:H_s>Tr_s\},
\]
then
\[
  Tr_{s_\star}
  \le\inf_s\{H_s+Tr_s\}.
\]
If a procedure cannot stop before $s_\star$ and returns $\tau\ge s_\star-1$, then
\[
  Tr_\tau
  \le2\inf_s\{H_s+Tr_s\}.
\]
\end{lemma}

\begin{proof}
If $s<s_\star$, then $Tr_s\ge2Tr_{s_\star}$.  If $s\ge s_\star$, then $H_s\ge H_{s_\star}>Tr_{s_\star}$.  The second conclusion follows from $r_\tau\le r_{s_\star-1}=2r_{s_\star}$.
\end{proof}

\begin{proof}[Proof of \cref{thm:profile-oracle}]
Every active set is nonempty and compact: the empirical primary maximizer
belongs to $G_s$ and remains in $U_s$. Thus the algorithm is defined on
every sample path. Each accepted stage obeys
$64\sum_{j\le s}|A_j|n_j=\Gamma_s\le Tr_s\le T$, so the exploration
fits in the horizon. Since every candidate set is nonempty and
$n_s\ge r_s^{-2}$, the stopping test eventually fails.

On the confidence event, \cref{lem:profile-invariants} bounds every
realized candidate size by the profile, including the first rejected
stage. Thus the prospective proxy at stage $s$ is at most
$H_U(s;\nu,\delta)$. If no stage is completed, the first prospective
cost exceeds $T$ and hence $H_U(0;\nu,\delta)>T$. Monotonicity then
makes the minimum with $T$ in the stated bound equal to $T$, so the
fallback policy suffices.

Otherwise, let $s_\star$ be the first deterministic crossing
$H_U(s;\nu,\delta)>Tr_s$. The algorithm cannot stop earlier, so its
last completed stage satisfies $\tau\ge s_\star-1$. Exploration
contributes at most $Tr_\tau/64$ positive regret, and commitment
contributes at most $6Tr_\tau$. The crossing lemma gives the displayed
oracle bound. Combining this with the deterministic cap $T$ gives the
minimum with $T$. On the confidence failure, positive regret is at
most $T$, which contributes the additive $\delta T$ term.
\end{proof}

\begin{proof}[Proof of \cref{cor:power-upper}]
Insert $\Pcal_\nu(r_j)\le Zr_j^{-d}$ into $H_U$.  The dyadic sum of $r_j^{-(d+2)}$ is dominated by its finest scale, yielding
\[
  H_U(s;\nu,T^{-2})
  \le C_dZr_s^{-(d+2)}\log(cZT)
\]
for scales that can improve on the trivial bound.  Balance this with $Tr_s$.
\end{proof}

\section{Regular Metric Lower Bound}
\label[appendix]{app:metric-lower}

We first give an explicit family satisfying \cref{ass:regular},
then prove the lower bound for every space in that assumption.
For $d>0$, let $\Sigma_d=\{0,1\}^{\N}$ with ultrametric
\[
  \rho_d(u,v)=a_d^{k(u,v)},
  \qquad a_d=2^{-1/d},
\]
where $k(u,v)$ is the number of common initial coordinates.  Take $M\asymp Z$ disjoint copies scaled to diameter $1/16$, put distance $1/4$ between copies, and add an anchor at distance one.  For $d=0$, use $\Theta(Z)$ isolated tie points.  This gives global packing growth $Zr^{-d}$ and local packing growth $(R/r)^d$.

Now fix any space satisfying \cref{ass:regular}. For
$0<\eps\le r_0$, choose a maximum $4\eps$-packing $V\subset S$.
Then
\[
  |V|=N\ge c_-Z\eps^{-d}.
\]
On the null,
\[
  \mu_0^1\equiv1/2,
  \qquad
  \mu_0^2(a)=g,
  \qquad
  \mu_0^2(z)=0\quad(z\in S).
\]
For $v\in V$, let
\[
  \mu_v^1(x)=\frac12+(\eps-D(x,v))_+,
  \qquad
  \mu_v^2=\mu_0^2.
\]
The hidden balls $B(v,\eps)$ are disjoint, and $\lambda_\star(\nu_v)=g/\eps$.

\begin{lemma}[Metric costly-search lower bound]
\label[lemma]{lem:metric-costly}
There are constants $c_d,c'_d>0$, depending only on $d$ and $c_-$,
such that, if
\[
  0<\eps\le\min\{g,r_0\},\qquad g\le1/8,
  \qquad
  T\eps^{d+3}\le c_dgZ,
\]
then every policy satisfies
\[
  \max\left\{
  \E_0R_2(T),
  \max_{v\in V}\E_vR_1(T)
  \right\}
  \ge c'_d\eps T.
\]
\end{lemma}

\begin{proof}
Let $N_v(T)$ count pulls in $B(v,\eps)$ and $N_S(T)$ pulls in $S$.  Under the null, $R_2(T)=gN_S(T)$.  If this expectation is below $c'_d\eps T$, disjointness and averaging give a $v$ with
\[
  \E_0N_v(T)
  \le\frac{c'_d\eps T}{gN}.
\]
The adaptive Gaussian chain rule gives
\[
  \KL(P_0\Vert P_v)
  \le\frac{\eps^2}{2}\E_0N_v(T)
  \le C_dc'_d\frac{T\eps^{d+3}}{gZ}.
\]
Choose constants so this is small.  Markov plus Pinsker gives
$P_v(N_v(T)\le T/2)\ge1/2$.  Outside the hidden ball the primary gap is exactly $\eps$, so $\E_vR_1(T)\ge\eps T/4$.
\end{proof}

\begin{proof}[Proof of \cref{thm:regular-metric}]
For the upper bound, an $r$-separated subset contains at most one
point in any ball of radius strictly less than $r/2$. Therefore, for
every $0<r\le1$,
\[
 \Pack(\X,r)
 \le1+\Cov\!\left(S,\min\{r/3,r_0\}\right)
 \le C_{d,c_+,r_0}Zr^{-d}.
\]
The primary near-optimality profile is bounded by this global packing
number. Apply \cref{cor:power-upper}, absorbing the fixed regularity
factor into the constant.

For the lower bound, fix $g=1/8$ and set
\[
 \eps=c_1\min\left\{r_0,(Z/T)^{1/(d+3)}\right\},
\]
where $c_1\le1$ is sufficiently small, depending only on $d$ and
$c_-$. Then $\eps\le r_0\le1/16<g$ and
$T\eps^{d+3}\le c_1^{d+3}Z$. The packing lower bound and
\cref{lem:metric-costly} give simultaneous regret at least
$c_2\eps T$. Finally, for $u>0$ and $r_0\le1$,
\[
 \min\{r_0,u\}\ge r_0\min\{1,u\}.
\]
This yields the stated rate, including the linear cap, with the
permitted dependence on $r_0$. The constructed null and alternatives
all have finite local trade-off, so the same lower bound holds on
that restricted class.
\end{proof}

\section{Arbitrary Entropy Proofs}
\label[appendix]{app:entropy}

Throughout this appendix, $\X=\{a\}\cup S$, $D(a,S)=1$,
$\operatorname{diam}(S)\le1/4$, and
$r_s=2^{-s-4}$. Write $M_s=\Pack(S,r_s)$ and
$p_s=1+M_s\le2M_s$. Every $r_s$-separated candidate net in $\X$
has at most $p_s$ points. The algorithms below use the realized nets;
neither algorithm is given $(M_s)_s$.

\subsection{Metric filtering with a supplied envelope}

For the unrestricted upper bound, run primary filtering as in
Algorithm~\ref{alg:pf-etc}, starting at radius $1/16$ instead of $1$.
For the supplied-envelope bound, set $B=1\vee L$ and
$q_s=\min\{1,Br_s\}$. The following metric version of robust-score
filtering gives an explicit implementation for the envelope case.

Initialize $U_{-1}=\X$ and $\Gamma_{-1}=0$. At stage $s$, form a
maximal $r_s$-separated subset $A_s$ of $U_{s-1}$ and set
\[
 n_s=\left\lceil 2r_s^{-2}\log\!\left(
 \frac{8\pi^2(s+1)^2\max\{2,|A_s|\}}{3\delta}\right)\right\rceil,
 \qquad
 \Gamma_s^{\mathrm{pro}}
 =\Gamma_{s-1}+64q_s|A_s|n_s.
\]
If $\Gamma_s^{\mathrm{pro}}>Tr_s$, stop before taking any stage-$s$
samples. Otherwise take $n_s$ fresh observations at each center, form
\[
 \widehat s_{B,s}(x)=2B\widehat\mu_s^1(x)+\widehat\mu_s^2(x),
\]
and retain
\[
 G_s=\left\{x\in A_s:
 \widehat s_{B,s}(x)\ge
 \max_{y\in A_s}\widehat s_{B,s}(y)-12Br_s\right\}.
\]
Set
\[
 U_s=U_{s-1}\cap\bigcup_{x\in G_s}\overline B(x,r_s),
 \qquad \Gamma_s=\Gamma_s^{\mathrm{pro}}.
\]
After the last completed stage $\tau$, commit to a maximizer of
$\widehat\mu_\tau^2$ over $G_\tau$. If no stage is completed, use any
fixed arm. Each active set is nonempty and compact, so the oracle
continues to be applicable even on a confidence failure.

\begin{lemma}[Metric robust-score invariants]
\label[lemma]{lem:metric-score-invariants}
On the simultaneous coordinate-confidence event, a representative
within distance $r_s$ of $x^\star$ is retained at every completed stage.
In particular, $x^\star\in U_s$. Moreover,
\[
 G_B(x)\le21Br_s\quad(x\in G_s),
 \qquad
 G_B(x)\le24Br_s\quad(x\in U_s).
\]
Every sampled stage-$s$ arm has maximum positive coordinate gap at
most $48q_s$. The terminal arm has
\[
 \Delta_1(\widehat x_\tau)\le21r_\tau,
 \qquad
 \Delta_2(\widehat x_\tau)\le3r_\tau.
\]
The number of exploration rounds is at most $T/64$.
\end{lemma}

\begin{proof}
The score $s_B=2B\mu^1+\mu^2$ is $(2B+1)$-Lipschitz, and its
empirical error is at most $(2B+1)r_s\le3Br_s$. If
$x^\star\in U_{s-1}$, maximality gives a center $\bar x_s$ at distance
at most $r_s$ from $x^\star$. Its true score gap is at most $3Br_s$,
so
\[
 \widehat s_{B,s}(\bar x_s)\ge s_B(x^\star)-6Br_s,
 \qquad
 \max_{x\in A_s}\widehat s_{B,s}(x)
 \le s_B(x^\star)+3Br_s.
\]
Thus $\bar x_s\in G_s$, which proves preservation of $x^\star$.
Every retained center has empirical score at least
$s_B(x^\star)-18Br_s$, hence true score gap at most $21Br_s$.
Expanding its ball by radius $r_s$ increases this bound by at most
$3Br_s$.

At stage $s\ge1$, the preceding expanded-set bound and
$r_{s-1}=2r_s$ give $G_B(x)\le48Br_s$. The robust-score sandwich
implies $\Delta_1(x)\le48r_s$ and
$[\Delta_2(x)]_+\le48Br_s$. Bounded means therefore give a maximum
positive gap of at most $48q_s$. At stage zero, $q_0\ge1/16$, so the
same conclusion follows from the unit bound on the gaps. At the
terminal stage, the score sandwich gives the primary conclusion,
and comparison with $\bar x_\tau$ gives
\[
 \mu^2(\widehat x_\tau)
 \ge\mu^2(\bar x_\tau)-2r_\tau
 \ge\mu^2(x^\star)-3r_\tau.
\]
Finally, $q_j\ge q_\tau\ge r_\tau$ for $j\le\tau$, and the completed
stage condition yields
\[
 \sum_{j\le\tau}|A_j|n_j
 \le\frac{\Gamma_\tau}{64q_\tau}
 \le\frac{Tr_\tau}{64q_\tau}\le\frac T{64}.
\]
This feasibility calculation holds on every sample path.
\end{proof}

\subsection{Entropy crossing}

Use $\delta=T^{-2}$ for $T\ge2$ and define
\[
 \bar n_s=\left\lceil2r_s^{-2}\log\!\left(
 \frac{8\pi^2(s+1)^2\max\{2,p_s\}}{3\delta}\right)\right\rceil.
\]
For $q_s^\infty=1$ and $q_s^B=\min\{1,Br_s\}$, put
\[
 H_s^\alpha=64\sum_{j=0}^s q_j^\alpha p_j\bar n_j,
 \qquad
 s_*^\alpha=\min\{s:H_s^\alpha>Tr_s\},
 \qquad \alpha\in\{\infty,B\}.
\]
These deterministic costs dominate their realized counterparts.
Consequently, neither algorithm can stop before $s_*^\alpha$.
When a stage is completed, its exploration and commitment regrets
are $O(Tr_\tau)$ on the confidence event, by the primary-filtering or
robust-score invariants. The last completed index obeys
$\tau\ge s_*^\alpha-1$. If no stage is completed, then
$s_*^\alpha=0$ and the fallback regret $T$ equals $16Tr_0$.
In all cases,
\begin{equation}
\label{eq:entropy-crossing-risk}
 \max_i\E R_i(T)\le CTr_{s_*^\alpha}+\delta T.
\end{equation}
The confidence proof is unchanged from
\cref{lem:stage-confidence}.

The first crossing lies in the range used by the critical entropy
functional. Indeed, $p_s\ge1$, $\bar n_s\ge r_s^{-2}$ and
$q_s^\alpha\ge r_s$, so $H_s^\alpha\ge64/r_s$. For
\[
 s_T=\left\lceil\log_2(16T)\right\rceil+2
\]
we have $r_{s_T}\le(1024T)^{-1}$ and hence
$H_{s_T}^\alpha>Tr_{s_T}$. Thus $s_*^\alpha\le s_T=O(\log(eT))$.

\begin{lemma}[Entropy crossing, unrestricted]
\label[lemma]{lem:entropy-cross-inf}
The unrestricted metric filtering algorithm satisfies
\[
 \max_i\E R_i(T)
 \le C\log^{2/3}(eT)\Psi_T^\infty(S).
\]
\end{lemma}

\begin{proof}
Write $s_*=s_*^\infty$ and $r_*=r_{s_*}$. At the first crossing,
some $j\le s_*$ satisfies
\[
 M_jr_j^{-2}\ell_j\ge\frac{cTr_*}{s_*+1},
 \qquad
 \ell_j=\log\bigl(cT^2(j+1)^2M_j\bigr),
\]
where universal constants absorb ceilings and $p_j\le2M_j$.
If $M_j\ge T$, the scale-$j$ term of $\Psi_T^\infty/T$ is
$r_j\ge r_*$. Otherwise, $\ell_j(s_*+1)\le C\log^2(eT)$, and
\[
 \frac{M_j}{T}
 \ge\frac{cr_*r_j^2}{\log^2(eT)}
 \ge\frac{cr_*^3}{\log^2(eT)}.
\]
Both $r_j$ and $(M_j/T)^{1/3}$ are therefore at least
$cr_*/\log^{2/3}(eT)$. Apply
\eqref{eq:entropy-crossing-risk}. The additive $\delta T$ is absorbed
because $M_0\ge1$ implies
$\Psi_T^\infty\ge\min\{T/16,T^{2/3}\}$.
\end{proof}

\begin{lemma}[Entropy crossing, supplied envelope]
\label[lemma]{lem:entropy-cross-B}
The metric robust-score algorithm satisfies
\[
 \max_i\E R_i(T)
 \le C\log(eT)\Psi_T^B(S).
\]
\end{lemma}

\begin{proof}
Write $s_*=s_*^B$ and $r_*=r_{s_*}$. Some $j\le s_*$ satisfies
\[
 q_j^B M_jr_j^{-2}\ell_j
 \ge\frac{cTr_*}{s_*+1}.
\]
If $M_j\ge T$, both $\sqrt{BM_j/T}$ and $(M_j/T)^{1/3}$ are at
least one, so the scale-$j$ contribution is $Tr_j\ge Tr_*$.
Otherwise set $F=C\log^2(eT)$ and obtain
\[
 q_j^B\frac{M_j}{T}\ge\frac{cr_*r_j^2}{F}.
\]
If $q_j^B=1$, then $Br_j\ge1$, and
\[
 \sqrt{\frac{BM_j}{T}}
 \ge c\sqrt{\frac{r_*r_j}{F}}\ge\frac{cr_*}{\sqrt F},
 \qquad
 \left(\frac{M_j}{T}\right)^{1/3}
 \ge c\left(\frac{r_*r_j^2}{F}\right)^{1/3}
 \ge\frac{cr_*}{F^{1/3}}.
\]
If $q_j^B=Br_j<1$, then
\[
 \frac{M_j}{T}\ge\frac{cr_*r_j}{BF}.
\]
The square-root term is again at least $cr_*/\sqrt F$.
Since $1/B>r_j$, the cubic term is at least
$cr_*/F^{1/3}$. Finally, $r_j\ge r_*$ and $F\ge1$, so all three
terms in the critical functional are at least
$cr_*/\sqrt F$. Apply \eqref{eq:entropy-crossing-risk} and absorb
$\delta T$ using
$\Psi_T^B\ge\min\{T/16,\sqrt T,T^{2/3}\}$.
\end{proof}

\subsection{One-scale lower bounds}

\begin{lemma}[One-scale costly-search lower bound]
\label[lemma]{lem:one-scale}
Let $V\subset S$ be an $r$-separated set of cardinality $N$.
If $0<\eps\le\min\{r/4,g\}$, $g\le1/8$, and
\[
 T\eps^3\le cgN,
\]
then the anchor-versus-hidden-cone family has simultaneous minimax
regret at least $c'\eps T$.
\end{lemma}

\begin{proof}
Use the null and cone alternatives of \cref{lem:metric-costly}.
The open balls $B(v,\eps)$ are disjoint. If the null-secondary
regret is smaller than $c'\eps T$, one of these balls receives at
most $c'\eps T/(gN)$ null-expected pulls. The trajectory KL to its
alternative is at most $\eps^2/2$ times this number. Choosing the
constants sufficiently small, Markov's inequality and Pinsker's
inequality imply that under this alternative at least $T/2$ pulls
occur outside the ball with probability at least $1/2$. Each such
pull has primary gap $\eps$.
\end{proof}

\begin{lemma}[One-scale primary-only lower bound]
\label[lemma]{lem:one-scale-primary}
For every $r$-separated set of size $N\ge1$ and every
$0<\eps\le r/4$ such that $T\eps^2\le cN$, there is a finite
family of admissible Gaussian instances with identical secondary
means and $\lambda_\star=0$ on which every policy incurs expected
primary regret at least $c'\eps T$ on some instance.
\end{lemma}

\begin{proof}
For $N\ge2$, use the reference primary mean $1/2$ and one
alternative $1/2+(\eps-D(\cdot,v))_+$ per packing point.
Let $N_v$ count pulls in $B(v,\eps)$. The chain rule gives
$\KL(P_0\Vert P_v)\le\eps^2\E_0N_v/2$, and Pinsker's inequality
therefore gives
\[
 \E_vN_v\le\E_0N_v+
 \frac{T\eps}{2}\sqrt{\E_0N_v}.
\]
Average over $v$ and use disjointness and Cauchy--Schwarz:
\[
 \frac1N\sum_v\E_vN_v
 \le\frac TN+\frac{T\eps}{2}\sqrt{\frac TN}.
\]
For a sufficiently small constant in $T\eps^2\le cN$, this is
at most $3T/4$. Some alternative therefore incurs primary regret
at least $\eps T/4$ outside its hidden ball.

For $N=1$, let $v$ be the packing point. In one instance the anchor
has primary mean $1/2+\eps$ and all of $S$ has mean $1/2$; in the
other, the anchor has mean $1/2$ and $S$ has mean
$1/2+(\eps-D(\cdot,v))_+$. These are Lipschitz instances, and their
trajectory KL is at most $T\eps^2/2$. If the anchor is played at
most $T/2$ times, the first instance incurs at least $\eps T/2$
regret; otherwise the second does. Bretagnolle--Huber gives a
constant lower bound on the sum of these two error probabilities.
\end{proof}

\begin{proof}[Proof of \cref{thm:entropy}]
The upper bounds are the two entropy-crossing lemmas.
For the unrestricted lower bound, choose a scale maximizing
$\Psi_T^\infty$ and an $r_s$-packing of cardinality $M_s$. Set
\[
 \eps=c\min\{r_s,(M_s/T)^{1/3}\}
\]
and use \cref{lem:one-scale} with $g=1/8$.

For the supplied-envelope lower bound, choose a scale maximizing
$\Psi_T^B$ and set
\[
 \eps=c\min\{r_s,\sqrt{BM_s/T},(M_s/T)^{1/3}\}.
\]
If $L\le1$, then $B=1$ and
$T\eps^2\le c^2M_s$, so \cref{lem:one-scale-primary} applies.
If $L\ge1$, set $g=\min\{L\eps,1/8\}$. A sufficiently small
universal $c$ ensures $\eps\le\min\{r_s/4,g\}$. In the
unsaturated case, the square-root cap gives
$T\eps^3\le c^2L\eps M_s=c^2gM_s$; in the saturated case,
the cubic cap gives $T\eps^3\le c^3M_s$. Thus
\cref{lem:one-scale} applies in both cases. Every alternative has
$\lambda_\star=g/\eps\le L$.
\end{proof}

\paragraph{Covering-tree implementation.}
A tree implementation requires a supplied hierarchy whose level-$s$
cells cover the relevant active set, have radius at most
$c_{\mathrm{cov}}r_s$, and have centers separated by at least
$c_{\mathrm{sep}}r_s$, with enumerable children. Expanding children
of retained cells implements the filtering construction after
adjusting its fixed geometric constants. Its non-sampling work is
linear in the number of generated nodes. If $W_s$ is the width of a
full level, packing and covering give
\[
 \Pack(S,3c_{\mathrm{cov}}r_s)
 \le W_s\le\Pack(S,c_{\mathrm{sep}}r_s).
\]
Consequently the width and packing descriptions agree up to fixed
radius rescaling. This statement assumes the hierarchy; it does not
supply a polynomial-time construction for an arbitrary black-box
metric.

\section{Pareto and High-Probability Proofs}
\label[appendix]{app:pareto}

\begin{lemma}[Sequential Bretagnolle--Huber]
\label[lemma]{lem:seq-bh}
Let $P,Q$ be laws of two adaptive bandit experiments. Include the policy randomization used to select $X_t$ in the prefix history, and let $A_t$ be measurable with respect to that history.  Then
\[
  P(A_t)+Q(A_t^c)
  \ge\frac12\exp\bigl(-\KL(P^{t-1}\Vert Q^{t-1})\bigr).
\]
\end{lemma}

\begin{proof}
Apply the Bretagnolle--Huber testing inequality to the two prefix-history laws and the test $A_t$. The common action-randomization kernel contributes zero to the relative entropy.
\end{proof}

\begin{proof}[Proof of \cref{thm:pareto}]
For alternative $j$, apply \cref{lem:seq-bh} to $A_t=\{X_t=j\}$.  The prefix KL is bounded by
\[
  \KL(P_0^T\Vert P_j^{\eps,T})
  =\frac{\eps^2}{2}\E_0N_j(T).
\]
Summing over $t$ yields
\[
  \E_0N_j(T)
  +\E_j^\eps(T-N_j(T))
  \ge\frac T2
  \exp\left(-\frac{\eps^2}{2}\E_0N_j(T)\right).
\]
Average over $j$ and use convexity of $e^{-x}$:
\[
  \frac1N\sum_j\E_0N_j(T)
  +\frac1N\sum_j\E_j^\eps(T-N_j(T))
  \ge\frac T2
  \exp\left(-\frac{\eps^2}{2N}\sum_j\E_0N_j(T)\right).
\]
Substitute $\sum_j\E_0N_j=r_0/g$ and upper-bound the average alternative regret by its maximum.  For the corollary, choose
\[
  \eps_0=
  \min\left\{
  \frac g2,
  \sqrt{\frac{gN}{4(r_0+gN)}}
  \right\}.
\]
Then the exponential is bounded below by a constant and $r_0/(gN)\le T/8$, so $r_\star(\eps_0)\ge c\eps_0T$.
\end{proof}

\section{High-Probability Lower and Upper Bounds}
\label[appendix]{app:high-prob}

\begin{lemma}[Strict binary-KL inversion]
\label[lemma]{lem:binary-kl}
Let $0<\delta\le1/16$.  If $p\ge9/16$ and
\[
 \kl(p,q)\le\frac18\log(1/\delta),
\]
then $q>\delta$.
\end{lemma}

\begin{proof}
Write $\ell=\log(1/\delta)$.  If $q\le\delta$, the binary entropy
bound gives
\[
 \kl(p,q)\ge p\log(1/q)-\log2
 \ge\frac9{16}\ell-\log2
 \ge\frac5{16}\ell>\frac18\ell,
\]
since $\ell\ge4\log2$.  This is a contradiction.
\end{proof}

\begin{proof}[Proof of \cref{thm:hp-lower}]
Use the Gaussian anchor-and-candidates family from
\cref{app:finite-lower}, with $g=1/4$, and put
$\ell=\log(1/\delta)$.  We prove that every policy incurs the claimed
regret with probability strictly greater than $\delta$ on at least one
instance.  Strictness is needed because the definition of
$\mathfrak Q_{T,K}(\delta)$ allows a tail probability equal to $\delta$.

\paragraph{Sublinear regime.}
Suppose
\[
 \eps=\left(\frac{N\ell}{T}\right)^{1/3}\le\frac18,
 \qquad
 m=\left\lfloor\frac{\ell}{16\eps^2}\right\rfloor.
\]
Then $m\ge\ell/(32\eps^2)\ge1$, while
\[
 \frac mT\le\frac{\ell}{16T\eps^2}
 =\frac{\eps}{16N}\le\frac1{128}.
\]
Let $N_S(T)=\sum_{j=1}^N N_j(T)$.  If
\[
 P_0\left(N_S(T)\ge\frac{Nm}{4}\right)>\delta,
\]
then with probability greater than $\delta$ the null-secondary regret
satisfies
\[
 R_2(T)=gN_S(T)\ge\frac{gNm}{4}
 \ge\frac{gN\ell}{128\eps^2}
 =\frac{g\eps T}{128}.
\]
Otherwise the event $A=\{N_S(T)<Nm/4\}$ has probability at least
$1-\delta$.  On every path in $A$, more than three quarters of the
candidates have been pulled fewer than $m$ times.  Averaging over the
candidates therefore gives some $j$ such that
\[
 p=P_0(E_j)\ge\frac34(1-\delta)\ge\frac9{16},
 \qquad E_j=\{N_j(T)<m\}.
\]
Stop the observation history at the $m$-th pull of arm $j$, or at time
$T$ if that pull never occurs.  The event $E_j$ is measurable with
respect to this stopped history.  Only observations from arm $j$ differ
between the null and alternative $j$, so the chain rule bounds the
stopped-history KL by
\[
 \KL(P_0^{\mathrm{stop}}\Vert P_j^{\mathrm{stop}})
 \le\frac{m\eps^2}{2}\le\frac{\ell}{32}.
\]
The same identity holds for randomized policies because their action
kernels are identical under both experiments.  By data processing and
\cref{lem:binary-kl}, $P_j(E_j)>\delta$.  On this event,
\[
 R_1(T)=\eps(T-N_j(T))\ge\eps(T-m)\ge\frac{\eps T}{2}.
\]
This proves the sublinear lower bound.

\paragraph{Linear regime.}
Suppose $(N\ell/T)^{1/3}>1/8$.  Set
\[
 \eps=\frac1{16},\qquad
 m=\max\left\{1,\left\lfloor\frac{T}{4096N}\right\rfloor\right\}.
\]
For $T\ge2$, one has $m\le T/2$.  Also $Nm\ge T/8192$: if
$T/(4096N)\ge2$, the floor is at least half its argument; otherwise
$N>T/8192$ and $m\ge1$.
Apply the same dichotomy as above.  If the null exploration event has
probability greater than $\delta$, then on that event
\[
 R_2(T)\ge\frac{gNm}{4}\ge\frac{T}{131072}.
\]
Otherwise some candidate $j$ has $P_0(E_j)\ge9/16$.  The stopped KL
satisfies
\[
 \KL(P_0^{\mathrm{stop}}\Vert P_j^{\mathrm{stop}})
 \le\frac m{512}
 \le\frac1{512}+\frac{T}{2097152N}
 <\frac{\ell}{8}.
\]
The last step uses $\ell>T/(512N)$ and $\ell\ge\log16>1$.
Again $P_j(E_j)>\delta$, and on $E_j$ the primary regret is at least
$\eps(T-m)\ge T/32$.

In each regime, choose the claimed regret threshold strictly below the
corresponding event's regret lower bound.  This proves the lower bound
for the strict-tail quantile $\mathfrak Q_{T,K}(\delta)$.
\end{proof}

\begin{lemma}[High-probability explore-filter-commit bound]
\label[lemma]{lem:hp-etc}
For $K\ge2$, $T\ge2$, and $0<\delta\le1/16$, there is a policy such
that
\[
 \Pp_\nu\left(
 \max_{i\le2}R_i^+(T)
 >C\min\{T,(K\log(4K/\delta))^{1/3}T^{2/3}\}
 \right)\le\delta
\]
uniformly over $\nu\in\F_K$.
\end{lemma}

\begin{proof}
Let $\ell=\log(4K/\delta)$ and $A=(K\ell)^{1/3}T^{2/3}$.
If $8A>T$, use an arbitrary arm.  Its regret is at most $T$, which is
within a universal factor of $\min\{T,A\}$ in this branch.
Otherwise set
\[
 n=\left\lceil\left(\frac{T^2\ell}{K^2}\right)^{1/3}\right\rceil,
 \qquad \beta=\sqrt{2\ell/n},
\]
sample each arm $n$ times, and apply the empirical primary filter and
secondary maximization in Algorithm~\ref{alg:lex-etc}.
The exploration is feasible: $8A\le T$ implies $K\ell\le T/512$,
so
\[
 Kn\le A+K\le T/8+T/512<T/2.
\]
A union bound gives
\[
 \Pp\left(\exists a,i:
 |\widehat\mu^i(a)-\mu^i(a)|>\beta\right)
 \le4K\exp(-n\beta^2/2)=\delta.
\]
On its complement, \cref{lem:finite-terminal} yields
\[
 \max_{i\le2}R_i^+(T)
 \le Kn+4T\beta\le C(K\ell)^{1/3}T^{2/3}.
\]
Together with the deterministic cap $T$, this proves the claim.
\end{proof}

\begin{proof}[Proof of \cref{thm:two-hp}]
Apply \cref{thm:hp-lower} with $N=1$ for the lower bound and
\cref{lem:hp-etc} with $K=2$ for the upper bound.  The upper confidence
factor is then $\log(8/\delta)$, as stated.
\end{proof}

\begin{proof}[Proof of \cref{prop:k-hp}]
Combine \cref{thm:hp-lower,lem:hp-etc} in the confidence and horizon
range of \cref{thm:hp-lower}.
\end{proof}

\section{Coefficient-Aware Comparison with Positive-Gap Annuli}

We record the annular coefficients for the explicit ultrametric family
in \cref{app:metric-lower}, with $\eps\le1/64$. This family has the
local covering property
\[
 \Cov(B(v,R)\cap S,r)\le C_d(R/r)^d,
 \qquad 0<r\le R<1/8,
\]
within the component containing $v$. This additional property is
needed below; global packing and covering bounds alone do not imply it.  In the notation of \citet{xue2025jmlr}, for a supplied valid parameter $\lambda$,
\[
  \X^i(r)
  =\left\{x:
  \frac{\Lambda_i(\lambda)r}{2}<\Delta_i(x)
  \le\Lambda_i(\lambda)r
  \right\}.
\]
Use the valid envelope $L=1$ on the null and the exact envelope on each alternative. On the null, all primary annuli are empty and a nonempty secondary annulus occurs only at $r\asymp g$, with covering coefficient $O_d(Zg^{-d})$.  On alternative $v$, positive secondary annuli are empty; primary annuli at $r<\eps$ lie in $B(v,O(r))$ and have $O_d(1)$ cover, while the $r\asymp\eps$ annulus has coefficient $O_d(Z\eps^{-d})$.  Thus all annular dimensions are zero.

With fixed $g$ and $\eps=(Z/T)^{1/(d+3)}\le1/64$, the known-trade-off annular benchmarks are, up to logarithms,
\[
  \sqrt{Zg^{-d}T}
  \quad\text{and}\quad
  \sqrt{Z\eps^{-d}T},
\]
whereas the adaptation lower bound is $\eps T$.  Their ratios are
\[
  \frac{\eps T}{\sqrt{Zg^{-d}T}}
  \asymp_d
  (T/Z)^{\frac{d+1}{2(d+3)}},
  \qquad
  \frac{\eps T}{\sqrt{Z\eps^{-d}T}}
  \asymp_d
  (T/Z)^{\frac{1}{2(d+3)}}.
\]
Both are polynomial when $T/Z$ grows polynomially.  Hence retaining the annular coefficient does not remove the separation.



\begin{thebibliography}{99}

\bibitem[Bubeck and Cesa-Bianchi(2012)]{bubeck2012survey}
S. Bubeck and N. Cesa-Bianchi.
\newblock Regret analysis of stochastic and nonstochastic multi-armed bandit problems.
\newblock \emph{Foundations and Trends in Machine Learning}, 5(1):1--122, 2012.

\bibitem[Bubeck et~al.(2011)Bubeck, Munos, Stoltz, and Szepesv\'ari]{bubeck2011xarmed}
S. Bubeck, R. Munos, G. Stoltz, and C. Szepesv\'ari.
\newblock $X$-armed bandits.
\newblock \emph{Journal of Machine Learning Research}, 12:1655--1695, 2011.
\newblock \href{https://www.jmlr.org/papers/v12/bubeck11a.html}{JMLR 12}.

\bibitem[Hadiji(2019)]{hadiji2019polynomial}
H. Hadiji.
\newblock Polynomial cost of adaptation for $X$-armed bandits.
\newblock In \emph{Advances in Neural Information Processing Systems 32}, pages 1029--1038, 2019.

\bibitem[H\"uy\"uk and Tekin(2021)]{huyuk2021lexicographic}
A. H\"uy\"uk and C. Tekin.
\newblock Multi-objective multi-armed bandit with lexicographically ordered and satisficing objectives.
\newblock \emph{Machine Learning}, 110:1233--1266, 2021.
\newblock \href{https://doi.org/10.1007/s10994-021-05956-1}{10.1007/s10994-021-05956-1}.

\bibitem[Kleinberg et~al.(2008)Kleinberg, Slivkins, and Upfal]{kleinberg2008metric}
R. Kleinberg, A. Slivkins, and E. Upfal.
\newblock Multi-armed bandits in metric spaces.
\newblock In \emph{Proceedings of the 40th Annual ACM Symposium on Theory of Computing}, pages 681--690, 2008.

\bibitem[Lattimore and Szepesv\'ari(2020)]{lattimore2020bandit}
T. Lattimore and C. Szepesv\'ari.
\newblock \emph{Bandit Algorithms}.
\newblock Cambridge University Press, 2020.

\bibitem[Locatelli and Carpentier(2018)]{locatelli2018adaptivity}
A. Locatelli and A. Carpentier.
\newblock Adaptivity to smoothness in $X$-armed bandits.
\newblock In \emph{Proceedings of the 31st Conference on Learning Theory}, PMLR 75:1463--1492, 2018.
\newblock \href{https://proceedings.mlr.press/v75/locatelli18a.html}{PMLR 75}.

\bibitem[Magureanu et~al.(2014)Magureanu, Combes, and Proutiere]{magureanu2014lipschitz}
S. Magureanu, R. Combes, and A. Proutiere.
\newblock Lipschitz bandits: Regret lower bound and optimal algorithms.
\newblock In \emph{Proceedings of the 27th Conference on Learning Theory}, PMLR 35:975--999, 2014.
\newblock \href{https://proceedings.mlr.press/v35/magureanu14.html}{PMLR 35}.

\bibitem[Skalse et~al.(2022)Skalse, Hammond, Griffin, and Abate]{skalse2022lexicographic}
J. Skalse, L. Hammond, C. Griffin, and A. Abate.
\newblock Lexicographic multi-objective reinforcement learning.
\newblock In \emph{Proceedings of the 31st International Joint Conference on Artificial Intelligence}, pages 3430--3436, 2022.

\bibitem[Slivkins(2011)]{slivkins2011implicit}
A. Slivkins.
\newblock Multi-armed bandits on implicit metric spaces.
\newblock In \emph{Advances in Neural Information Processing Systems 24}, pages 1602--1610, 2011.

\bibitem[Tekin and Turgay(2018)]{tekin2018dominant}
C. Tekin and E. Turgay.
\newblock Multi-objective contextual multi-armed bandit with a dominant objective.
\newblock \emph{IEEE Transactions on Signal Processing}, 66(14):3799--3813, 2018.

\bibitem[Turgay et~al.(2018)Turgay, Oner, and Tekin]{turgay2018similarity}
E. Turgay, D. Oner, and C. Tekin.
\newblock Multi-objective contextual bandit problem with similarity information.
\newblock In \emph{Proceedings of the 21st International Conference on Artificial Intelligence and Statistics}, PMLR 84:1673--1681, 2018.

\bibitem[Xue et~al.(2024)Xue, Cheng, Liu, Wang, and Zhang]{xue2024lipschitz}
B. Xue, J. Cheng, F. Liu, Y. Wang, and Q. Zhang.
\newblock Multiobjective Lipschitz bandits under lexicographic ordering.
\newblock In \emph{Proceedings of the 38th AAAI Conference on Artificial Intelligence}, 38(15):16238--16246, 2024.
\newblock \href{https://doi.org/10.1609/aaai.v38i15.29558}{10.1609/aaai.v38i15.29558}.

\bibitem[Xue et~al.(2025a)Xue, Cheng, Liu, Wang, Zhang, and Zhang]{xue2025jmlr}
B. Xue, J. Cheng, F. Liu, Y. Wang, L. Zhang, and Q. Zhang.
\newblock Lexicographic Lipschitz bandits: New algorithms and a lower bound.
\newblock \emph{Journal of Machine Learning Research}, 26(223):1--56, 2025.
\newblock \href{https://www.jmlr.org/papers/v26/24-1047.html}{JMLR 26(223)}.

\bibitem[Xue et~al.(2025b)Xue, Lin, Zhang, and Zhang]{xue2025linear}
B. Xue, X. Lin, X. Zhang, and Q. Zhang.
\newblock Multiple trade-offs: An improved approach for lexicographic linear bandits.
\newblock In \emph{Proceedings of the 39th AAAI Conference on Artificial Intelligence}, 39(20):21850--21858, 2025.
\newblock \href{https://doi.org/10.1609/aaai.v39i20.35491}{10.1609/aaai.v39i20.35491}.

\bibitem[Xue et~al.(2025c)Xue, Lin, Wan, and Zhang]{xue2025corruptions}
B. Xue, X. Lin, Y. Wan, and Q. Zhang.
\newblock Problem-dependent regret for lexicographic multi-armed bandits with adversarial corruptions.
\newblock In \emph{Proceedings of the 34th International Joint Conference on Artificial Intelligence}, pages 6776--6784, 2025.
\newblock \href{https://doi.org/10.24963/ijcai.2025/754}{10.24963/ijcai.2025/754}.

\bibitem[Xue et~al.(2026)Xue, Wan, Lu, and Zhang]{xue2026lexelim}
B. Xue, Y. Wan, Z. Lu, and Q. Zhang.
\newblock Beyond the lower bound: Bridging regret minimization and best arm identification in lexicographic bandits.
\newblock In \emph{Proceedings of the 40th AAAI Conference on Artificial Intelligence}, 40(32):27414--27422, 2026.
\newblock \href{https://doi.org/10.1609/aaai.v40i32.39959}{10.1609/aaai.v40i32.39959}.

\end{thebibliography}

\end{document}
